\documentclass[11pt]{amsart}

\usepackage[T1]{fontenc}
\usepackage[utf8]{inputenc}
\usepackage{microtype}
\usepackage{mathtools}
\usepackage{amssymb}
\usepackage{comment}
\usepackage{enumitem}
% Shared review markup and formal links; the human draft hides declaration names.
% Review annotations, following the red-add/blue-remove convention used in
% the Varenna 2016 papers. These commands keep the original wording visible.
% Accept an insertion by removing its AIaddition environment or AIadd wrapper.
% Accept a replacement by deleting AIremove{old} and unwrapping AIadd{new}.
% Resolve a comment by deleting its complete AIcomment{ID}{text} command.
% There is deliberately no global switch that hides unresolved comments.
\usepackage{xcolor}
\usepackage{xurl}
\usepackage[normalem]{ulem}
% Give long inline review markup a little extra line-breaking room without
% changing the author text or the mathematical content.
\setlength{\emergencystretch}{.75em}
% AI compile repair: the retained draft uses the undefined typo \cots.
% EDIT-02 beside that occurrence makes this repair visible in the PDF.
\providecommand{\cots}{\cdots}
% Permit a line break after commas in long inline coordinate lists.
\begingroup\lccode`~=`,
\lowercase{\endgroup\def~}{\mathchar"613B\allowbreak}
\mathcode`,="8000
\definecolor{AIadded}{RGB}{160,25,25}
\definecolor{AIreview}{RGB}{25,70,145}
\newcommand{\AIadd}[1]{{\color{AIadded}#1}}
\newcommand{\AIremove}[1]{{\color{AIreview}%
  \ifmmode\text{\sout{\ensuremath{#1}}}\else\sout{#1}\fi}}
\newcommand{\AIreplace}[2]{\AIremove{#1}\,\AIadd{#2}}
\newenvironment{AIaddition}[1]
  {\par\begingroup\color{AIadded}\noindent
   \textbf{AI addition: #1.}\par\smallskip\ignorespaces}
  {\par\endgroup}
\newcommand{\AIcomment}[2]{%
  \par\smallskip\noindent
  \fcolorbox{AIreview}{AIreview!4}{%
    \parbox{\dimexpr\linewidth-2\fboxsep-2\fboxrule\relax}{%
      \normalfont\small\color{AIreview}\textbf{AI comment #1.} #2}}%
  \par\smallskip}
\newcommand{\AIreviewlegend}{%
  \AIcomment{Review guide}{Red text proposes additions or corrections;
  blue struck text preserves the original wording.
  \textbf{Johanna:} review the proposals as mathematics and prose, and resolve
  their boxes after accepting or revising them. No Lean installation, code
  editing, or proof replay is required for her review.
  \textbf{Max:} review the whole paper--Lean correspondence, including
  every statement, its hypotheses, and the definitions entering the challenge. The technical supplement contains the comparison notes, mathematical accounts,
  and optional declaration links. The
  side-by-side review interface links statements to their checked sources.}}

% Lean cross-references for the paper--Lean audit.
% Each formalized result is followed by \formalref{...}; every \leandecl or
% \leanlib names one declaration by file, line and fully qualified name at
% the pinned revisions below. The review generator in itpplasma/stafford38-formal
% (tools/paper_lean_audit) checks every such reference against the Lean
% source, so a stale file, line or name fails its build.
% Accepting: keep this file and the \formalref lines, or delete both.
% Pins: stafford38-formal v1.3.1 (current Lean source; DOI 10.5281/zenodo.23127367),
% AlgebraicAnalysis v0.3.3, global-stafford-formal v1.0.6, and
% symplectic-stafford-formal main (unreleased).
\newcommand{\leanrepo}{https://github.com/itpplasma/stafford38-formal/blob/f9448d6307fa25aff9d9f94ce0a9c7f9b03e63ff}
\newcommand{\leanlibrepo}{https://github.com/itpplasma/algebraic-analysis/blob/bbbbf3fc358ca8100b158cec4cf47f336ab70163}
\newcommand{\leanglobalrepo}{https://github.com/itpplasma/global-stafford-formal/blob/e676caed1253e5d09fc6699d5dc683febbabf7bb}
\newcommand{\leanfilteredrepo}{https://github.com/itpplasma/symplectic-stafford-formal/blob/dccff0d8fe723ec1478d1f63bbbdf648ba240bfa}
% \leandecl{file}{line}{Fully.Qualified.Name}: declaration in stafford38-formal.
\newcommand{\leandecl}[3]{\href{\leanrepo/#1\#L#2}{\nolinkurl{#3}}}
% \leanlib{file}{line}{Fully.Qualified.Name}: declaration in AlgebraicAnalysis.
\newcommand{\leanlib}[3]{\href{\leanlibrepo/#1\#L#2}{\nolinkurl{#3}}}
% \leanglobal, \leanfiltered: declarations in the two generalizations of Section 1.2.
\newcommand{\leanglobal}[3]{\href{\leanglobalrepo/#1\#L#2}{\nolinkurl{#3}}}
\newcommand{\leanfiltered}[3]{\href{\leanfilteredrepo/#1\#L#2}{\nolinkurl{#3}}}
% The human-readable paper displays only substantive mathematical qualifications.
% Full file paths and declaration names remain in source and in the supplement.
\newif\ifshowleandetails
\showleandetailsfalse
% \formalref[relation to the printed statement]{declarations}
\newcommand{\formalref}[2][]{%
  \ifshowleandetails
    \par\smallskip\noindent
    {\small\textsf{Lean}\if\relax\detokenize{#1}\relax\else\ (#1)\fi: #2}%
    \par\smallskip
  \else
    \if\relax\detokenize{#1}\relax\else
      \par\smallskip\noindent{\small\textsf{Formal comparison:} #1}%
      \par\smallskip
    \fi
  \fi}

\usepackage[colorlinks=true,citecolor=blue,linkcolor=blue,urlcolor=blue]{hyperref}
\numberwithin{equation}{section}

\newtheorem{theorem}{Theorem}[section]
\newtheorem{proposition}[theorem]{Proposition}
\newtheorem{lemma}[theorem]{Lemma}
\newtheorem{corollary}[theorem]{Corollary}
\theoremstyle{plain}
\newtheorem{definition}[theorem]{Definition}
\newtheorem{convention}[theorem]{Convention}
\theoremstyle{remark}
\newtheorem{remark}[theorem]{Remark}



\title{Stafford's conjecture on cyclicity of torsion modules}
\author{Christopher Albert, Johanna Moser and Maximilian Philipp}
\address{Graz University of Technology}
\date{}
\newcommand{\Char}{\operatorname{Char}}
\newcommand{\Weyl}{A_n}
\newcommand{\ord}{\operatorname{ord}}
\newcommand{\gr}{\operatorname{gr}}
\newcommand{\Ann}{\operatorname{Ann}}
\newcommand{\Dir}{\operatorname{Dir}}
\newcommand{\Sp}{\operatorname{Sp}}
\newcommand{\Spec}{\operatorname{Spec}}
\newcommand{\Supp}{\operatorname{Supp}}
\newcommand{\Sing}{\operatorname{Sing}}
\newcommand{\sm}{\mathrm{sm}}
\newcommand{\bbA}{\mathbb{A}}
\newcommand{\bbP}{\mathbb{P}}
\newcommand{\bbN}{\mathbb{N}}
\newcommand{\bbZ}{\mathbb{Z}}
\newcommand{\bbQ}{\mathbb{Q}}
\newcommand{\Kbar}{\overline{k}}
\newcommand{\dd}{\mathrm{d}}
\DeclareMathOperator{\rad}{rad}


\begin{document}
 \setlength{\parindent}{0pt}
\begin{abstract}
Let $A_n(k)$ be the Weyl algebra over a field of characteristic zero.  We prove
that every nonzero $d\in A_n(k)$ admits $F,R,S\in A_n(k)$ such that
$1=dR+FdS$.  Thus Stafford's Conjecture 3.8 holds.  \AIadd{For $n\ge1$,} \AIreplace{The}{the} multiplier $F$ may be
chosen as $x^N$, after a linear symplectic change of variables, where $N$ is
the Bernstein degree of $d$.  The proof studies a cyclic quotient through its
characteristic variety and excludes nonzero quotients by conormal geometry at
infinity.
\end{abstract}

\maketitle
\AIreviewlegend


\section{Introduction}

Let $k$ be a field of characteristic zero, and let
\[ A_n(k)=k\langle x_1,\ldots,x_n,\partial_1,\ldots,\partial_n\rangle\]
be the $n$th Weyl algebra, with $[\partial_i,x_j]=\delta_{ij}$ and
$[\partial_i,\partial_j]=[x_i,x_j]=0$.

\AIremove{Stafford conjectured in 1978 (Conjecture~3.8 of
\cite[p.~438]{Stafford1978}) that for every nonzero $d\in A_n(k)$, one can
find $F\in A_n(k)$ with $A_n(k)=dA_n(k)+FdA_n(k)$. Stafford notes that
this immediately implies $A_n(k)/dA_n(k)$ is cyclic, i.e.\ generated by a
single element. \cite{Bellamy2026} reported the resulting question --
whether every right torsion module over $A_n(k)$ is cyclic -- as open
[Section~3, Conjecture~3.4]. We prove Stafford's conjecture:}\par
\AIadd{Stafford's Conjecture~3.8 asks whether every nonzero $d\in A_n(k)$
admits $F\in A_n(k)$ with $A_n(k)=dA_n(k)+FdA_n(k)$
\cite[p.~438]{Stafford1978}. Stafford identifies this with stable cyclicity
for finitely generated torsion right Weyl-algebra modules
\cite[pp.~437--438]{Stafford1978}. Bellamy's survey
states the weaker ordinary-cyclicity question for finitely generated torsion
$A_n(\mathbb C)$-modules as open~\cite[Section~3, Conjecture~3.4]{Bellamy2026}.
We prove Stafford's conjecture:}
\AIcomment{INTRO-01}{\textbf{Proposed change --- author acceptance pending.} The quotient $A_n/dA_n$ is always cyclic, generated
by the class of $1$. The nontrivial implication concerns general
\emph{finitely generated} torsion modules. Bellamy states the ordinary
cyclicity question over $\mathbb C$; Stafford's stable-cyclicity formulation
is stronger. The red replacement above restores these distinctions.}

\begin{theorem}\label{thm:main}
For every characteristic-zero field $k$, every $n\geq1$, and every
nonzero $d\in A_n(k)$, there exist $F,R,S\in A_n(k)$ such that
\[ 1=dR+FdS.\]
Moreover, one may choose $F=\ell^{\deg_Bd}$ for a linear Weyl coordinate
$\ell$, obtained by an invertible linear symplectic change of generators.
\AIadd{Explicitly, $\ell=\Phi_g^{-1}(x_1)$ for a linear
$\omega$-preserving automorphism $g$ as in Section~\ref{sec:monic}.}
\end{theorem}
\AIadd{\formalref[{Lean also covers $n=0$, where $A_0(k)=k$}]{%
\leandecl{Stafford38/FoundationClosure.lean}{67}{Stafford38.universalStatement} proves
\leandecl{Stafford38/Statement.lean}{23}{Stafford38.UniversalStatement};
\leandecl{Stafford38/FoundationClosure.lean}{71}{Stafford38.universalFixedSourceStatement} proves
\leandecl{Stafford38/FixedSourceStatement.lean}{35}{Stafford38.FixedSource.UniversalFixedSourceStatement}
with \leandecl{Stafford38/FixedSourceStatement.lean}{28}{Stafford38.FixedSource.IsLinearWeylCoordinate}
and \leandecl{Stafford38/FixedSourceStatement.lean}{23}{Stafford38.FixedSource.bernsteinDegree};
Mathlib-only statements compared by Comparator:
\leandecl{Stafford38/ChallengeDefinitions.lean}{78}{Stafford38Challenge.UniversalStatement},
\leandecl{Stafford38/ChallengeDefinitions.lean}{193}{Stafford38FixedSourceChallenge.UniversalFixedSourceStatement}.}}

\AIadd{The machine-checked proof of this result is registered in
Palomar~\cite{Stafford38Palomar2026,Stafford38Zenodo123}. Its Global Stafford extension
\cite{GlobalStaffordPalomar2026,GlobalStaffordZenodo106} proves the same-divisor identity
$1=dR+FdS$ in $D_k(A)$ for every smooth integral finitely generated
characteristic-zero $k$-algebra $A$ and every nonzero finite-order
$k$-linear differential operator $d$. Polynomial $A$ recovers the
Weyl-algebra identity; the global theorem asserts existence of $F$,
without the fixed linear-coordinate refinement in Theorem~\ref{thm:main}.}

\textit{Proof sketch:}
Give each standard generator Bernstein degree one and
let $N=\deg_B d$.  A linear symplectic change of variables makes $d$ monic of
degree $N$ in a momentum $p$, whose conjugate coordinate we denote by $x$.  It is
enough to prove that the cyclic right module
\begin{equation} Q=A_n/(dA_n+x^NdA_n).
 \label{eq:Q-intro}\end{equation}
vanishes.  An Euler identity makes $x$ surjective on $Q$, while monicity makes
$H=\{x=0\}$ noncharacteristic and places $\Char(Q)$ in a hypersurface
$V(P)$ with $P(dx)\ne0$.  A filtered-page length argument shows that
$\Char(Q)$ has no point over $H$.  Gabber's theorem and the conormal
argument of \AIreplace{Sections 4 and 5}{Sections~\ref{sec:conormal-review}
and~\ref{sec:involutive-review}} then force $\Char(Q)$ to be empty.  The same
argument gives the more precise identity $A_n=dA_n+x^NdA_n$ in the chosen
symplectic coordinates.

\AIadd{Constructive predecessors concern Stafford's established generator
theorems. Hillebrand--Schmale and Leykin give effective two-generator
methods, and Leykin also constructs cyclic generators of holonomic
modules~\cite{HillebrandSchmale2001,Leykin2004}. Quadrat--Robertz develop
module-decomposition and basis algorithms~\cite{QuadratRobertz2007,
QuadratRobertz2014}; Caro--Levcovitz and Caro-Tuesta--Levcovitz study
differential-operator rings with series coefficients~\cite{CaroLevcovitz2010,
CaroTuestaLevcovitz2020}. The constrained identity for arbitrary nonzero $d$
requires the additional quotient-vanishing argument developed here.}

\AIadd{The surrounding generator theory includes Stafford's stable-range
and noncommutative Forster--Swan results~\cite{StaffordStableRange1981,
StaffordGenerating1981}. The generator article has a corrigendum and
related exposition~\cite{StaffordGeneratingCorrigendum1983,
StaffordGeneratingChapter1982}. Stafford also constructed stably free
Weyl-algebra ideals~\cite{StaffordStablyFree1985}.
OreModules supplies computational methods for modules over Ore algebras
\cite{OreModules2007}. These provide context; the Euler residue,
filtered-length contradiction, and asymptotic conormal construction below
are project arguments.}

\section{\texorpdfstring{\AIremove{AI Statement [Chris]}\AIadd{Formal verification}}{Formal verification}}
\AIreplace{Every statement of this paper, including
Theorems~\ref{thm:asymptotic-conormal} and \ref{thm:coisotropic-exclusion}
at the generality stated, is formalized and verified in Lean~4 against
Mathlib, with no axiom beyond Lean's three standard ones.}{The main identity, its fixed-source strengthening and the documented
cyclicity and noncharacteristic consequences are verified in Lean~4 against
Mathlib using only the three standard axioms. The formal proof also supplies
the characteristic-support and geometric results used by its own argument.
Statement scope and proof routes must be compared separately: machine
verification of a terminal theorem does not certify every intermediate
step of the printed proof. The mathematical differences are noted below;
full formal comparisons and declaration links are in the technical supplement.}
\AIadd{The headline theorem is registered in Palomar~\cite{Stafford38Palomar2026};
the current Lean sources and auxiliary results are archived~\cite{Stafford38Zenodo131}.}
The accompanying proof-map supplement names the Lean declaration for each step.
\AIcomment{WF-01}{\textbf{Proposed change --- author acceptance pending.} Please review each comparison at the stated generality.
The Palomar record checks the registered main theorem; later results have
separate archived sources. The formal geometric route uses Laurent series,
while the printed tangent-limit argument is auxiliary. The short comparison
notes identify these differences without claiming that every numbered
statement or every proof passage has been verified verbatim.}
\AIadd{The development uses Lean~4~\cite{Lean4}, Mathlib~\cite{Mathlib2020}
and AlgebraicAnalysis v0.3.0~\cite{AlgebraicAnalysis030}; tooling, chronology
and statement comparison are described in Appendix~\ref{app:ai-workflow}.}



















\AIadd{\paragraph{Online resources.} Online Resource~1 is the interactive paper--Lean correspondence review, including the exact challenge definitions and archived proof endpoints; Online Resource~2 is the readable mathematical account of the selected formal proofs. Both are available from the \href{https://itpplasma.github.io/stafford38-formal/}{review companion}, with a \href{https://github.com/itpplasma/stafford38-supplementary}{frozen supplementary package} for journal review~\cite{Stafford38Supplementary021}. The reusable HTML generator is maintained separately under the Apache~2.0 license; its software archive and the manuscript supplement have distinct version records. Proposed corrections and unresolved review actions remain explicitly marked in this draft.}

\section{Notation and definitions}\label{sec:setting}


Throughout, $k$ is a field of characteristic zero and for the $n$-th Weyl algebra over $k$, we write
\[ A_n(k)=k\langle x_1,\dots,x_n,\partial_1,\dots,\partial_n\rangle,
  \quad [\partial_i,x_j]=\delta_{ij},\quad [x_i,x_j]=[\partial_i,\partial_j]=0 .\]
All ideals are \emph{right} ideals and all modules are \emph{right} modules
unless the contrary is said explicitly.

\begin{convention}\label{conv:pair}
We single out the first Weyl pair $(x_1, \partial_1)$ and write
\[\qquad p:=\partial_1,\qquad
  C:=A_{n-1}(k)[x_1],\qquad A_n(k)=C[p;\delta],\]
Here $\delta(x_1)=1$ and $\delta$ vanishes on $A_{n-1}(k)$.
so that $x_1$ is central in $C$, $[p,x_1] = 1$, and the transverse algebra $A_{n-1}(k)$ commutes with $x_1$ and $p$.
Every element of $A_n(k)$ can uniquely be written as $\sum_j p^j c_j$ with $c_j\in C$.
\end{convention}
\AIadd{\formalref{%
\leandecl{Stafford38/Weyl/IteratedEquivalence.lean}{636}{Stafford38.WeylIteratedEquivalence.presentedIteratedEquiv},
\leandecl{Stafford38/Weyl/OuterOreMonic.lean}{47}{Stafford38.WeylOuterOreMonic.presentedOuterPolynomial}}}




\begin{definition}[Filtration {[Def.~2.6 and ~2.9 \cite{SchnellNotes}]}]\label{def:filtration}
A filtration $F_\bullet$ on $A_n(k)$ is a sequence of linear subspaces
\[\{0\} = F_{-1} \subseteq F_0 \subseteq F_1\subseteq...\subseteq A_n(k),\]
such that $F_i\cdot F_j \subseteq F_{i+j}$ and $A_n(k) = \bigcup F_j.$
The associated graded algebra, also called symbol ring, is defined as
\[\mathrm{gr}^F A_n(k) = \bigoplus_{j = 0}^\infty F_j/F_{j-1}.\]
Let $a\neq 0 \in A_n(k)$ and let $j$ be the smallest integer, such that $a \in F_j$, called a \emph{filtration degree} $\deg_F a$. The class of $a$ in $F_j/F_{j-1}$ is called its \emph{symbol} $\sigma_j(a)$.
\end{definition}

In the proof, we use two filtrations of $A_n(k)$.
\begin{itemize}[nosep,leftmargin=1.4em]
\item The \emph{Bernstein filtration} $B_\bullet$ assigns degree $1$ to every
  $x_i$ and every $\partial_i$. The Bernstein degree of $d$ is
  written $\deg_B d$.
\item The \emph{order filtration} $F_\bullet$ assigns degree $0$ to every $x_i$
  and degree $1$ to every $\partial_i$; $F_m$ is spanned by the monomials
  $x^\alpha\partial^\beta$ with $|\beta|\le m$.  %We write $\sigma_m(a)$ for the order-$m$ principal symbol.
\end{itemize}
\AIadd{\formalref[{as spans of ordered PBW words of bounded weight}]{%
\leandecl{Stafford38/Weyl/Filtration.lean}{284}{Stafford38.WeylFiltration.bernsteinPiece},
\leandecl{Stafford38/Weyl/Filtration.lean}{288}{Stafford38.WeylFiltration.orderPiece},
\leandecl{Stafford38/FixedSourceStatement.lean}{23}{Stafford38.FixedSource.bernsteinDegree}}}
For both the order filtration and the Bernstein filtration, the associated graded algebra $\mathrm{gr} A_n(k)$ is simply the polynomial ring $k[x,\xi]=k[x_1,\dots,x_n,\xi_1,\dots,\xi_n]$ in $2n$ variables with $\xi_i$ the symbol of $\partial_i$. In particular, $\mathrm{gr} A_n(k)$ is commutative \cite[\AIadd{Prop.~2.10}]{SchnellNotes}.
Good filtrations and characteristic varieties are treated systematically for example in \cite{Bjork,HTT}. %Bj\"ork and Hotta--Takeuchi--Tanisaki


\subsection{Cotangent space and symplectic structure}\label{sec:cotangent}


Let $X=\mathbb A^n_k$ be affine $n$-space over a field $k$ of characteristic zero, with coordinates
$x_1,\ldots,x_n$. Its \emph{cotangent bundle} $T^*X$ is the space of pairs
$(y,\xi)$ consisting of a point $y\in X$ together with a covector $\xi$
at $y$; since $X$ is affine space, this cotangent space is the same
vector space $(k^n)^*$ at every point, and we may identify
\[
 T^*X \;\cong\; X\times(k^n)^*.
\]
We write $\xi_1,\ldots,\xi_n$ for the coordinates on $(k^n)^*$ dual to
$x_1,\ldots,x_n$, so that a covector is $\sum_i\xi_i\,dx_i$, and identify
$\gr^BA_n(k)=k[x,\xi]$ with the coordinate ring of $T^*X$. We write
$\pi\colon T^*X\to X$ for the projection forgetting the covector, and for
$y\in X$ write $T_y^*X=\pi^{-1}(y)\cong(k^n)^*$ for the \emph{fibre} of
$T^*X$ over $y$.
\AIadd{\formalref{\leandecl{Stafford38/Characteristic/Polynomial.lean}{22}{Stafford38.Characteristic.SymbolRing}
(the ring $k[x,\xi]$; the first summand indexes base coordinates $x_i$, and the second indexes fibre coordinates $\xi_i$)}}

We single out the first coordinate $x_1$ and define the distinguished
hyperplane
\[
 H:=\{x_1=0\}=\{(0,x_2,\ldots,x_n)\}\subseteq X;
\]
its conormal covector is $dx_1$, i.e.\ $\xi=(1,0,\ldots,0)$.


\begin{definition}\label{def:good-filtration}
A filtration of a right $A_n$-module $M$ compatible with $F_\bullet$
is an exhaustive increasing chain of $k$-subspaces $M_\bullet$ with
$M_iF_j\subseteq M_{i+j}$ and $M_i=0$ for $i<0$. It is \emph{good} if the associated graded module
$\gr M=\bigoplus_iM_i/M_{i-1}$ is finitely generated over $\gr^FA_n$.
\end{definition}
As an example, the filtration $M_i:=qF_i$ of a cyclic module over $A_n$ induced by its generator $q$ is good, since $\gr Q$ is generated by the element $\bar{q}$ denoting the class of $q$ in $Q_0/Q_{-1}.$

For a finitely generated right $A_n$-module $M$ with a good order filtration
$M_\bullet$, the graded module $\gr M$ is a finitely generated $k[x,\xi]$-module.
\begin{definition}\label{def:char-variety}
    The \emph{characteristic variety} (also called symbol variety) of a finitely generated $A_n$-module $M$ with a good order filtration
$M_\bullet$ is defined by
\[
  \Char(M):=V\bigl(\Ann_{k[x,\xi]}\gr M\bigr)\subseteq
  T^*\bbA^n=\Spec k[x,\xi],
\]
$\Ann_{k[x,\xi]}\gr M = \{f \in k[x,\xi]: f\cdot m = 0\,\forall m \in \gr M\}$ are the \emph{annihilators} of $\gr M$ over $k[x,\xi]$, and $V$ denotes their \emph{vanishing locus}, i.e. the set of points $(x, \xi)$ for which all annihilator polynomials vanish. This set of points lies in the \emph{cotangent bundle} $T^*\bbA^n$ of the affine $n$-space $\bbA^n.$

$\Char(M)$ is independent of the choice of the good filtration.
\end{definition}
\AIadd{\formalref[{for $M=A_n/I$}]{%
\leandecl{Stafford38/Characteristic/InitialIdeal.lean}{189}{Stafford38.CharacteristicInitialIdeal.orderInitialIdeal},
\leandecl{Stafford38/Characteristic/InitialIdeal.lean}{197}{Stafford38.CharacteristicInitialIdeal.orderCharacteristicSupport},
\leandecl{Stafford38/Characteristic/AssociatedGradedModule.lean}{373}{Stafford38.CharacteristicAssociatedGradedModule.annihilator_orderAssociatedGradedModule}}}


\begin{definition}[Symplectic form and Poisson bracket]\label{def:symplectic-poisson}
A bilinear form $\omega$ on a $k$-vector space is \emph{symplectic} if it
is antisymmetric and nondegenerate.

The standard symplectic form on the vector space $k^{2n}$, with basis
$x_1,\ldots,x_n,\partial_1,\ldots,\partial_n$, is determined by the
relations
\[ \omega(x_i,\partial_j)=\delta_{ij},\qquad
 \omega(x_i,x_j)=\omega(\partial_i,\partial_j)=0,\]
analogously to the Weyl relations. Identifying $T^*X\cong X\times(k^n)^*$ via the coordinates $x_1,\ldots,x_n$ and their duals
$\xi_1,\ldots,\xi_n$, on $T^*X$ this $\omega$ takes the familiar form
\[ \omega=\sum_{i=1}^n dx_i\wedge d\xi_i.\]

The associated \emph{Poisson bracket} on $k[x,\xi]$, the coordinate ring
of $T^*X$, is
\[ \{f,g\}=\sum_{i=1}^n\left(
   \frac{\partial f}{\partial x_i}\frac{\partial g}{\partial \xi_i}
   -\frac{\partial f}{\partial \xi_i}\frac{\partial g}{\partial x_i}
 \right).\]
\end{definition}
With this convention of the Poisson bracket, the order-filtration symbol of a commutator is $\sigma^F([a,b])$ is $\{\sigma^F(b),\sigma^F(a)\}$.
\AIadd{\formalref[{}]{%
\leandecl{Stafford38/Characteristic/Polynomial.lean}{31}{Stafford38.Characteristic.poissonBracket},
\leandecl{Stafford38/Characteristic/SymplecticCompletion.lean}{31}{Stafford38.CharacteristicSymplecticCompletion.phasePairing}}}

\begin{definition}[involutive]\label{def:involutive}
An ideal $J$ is \emph{involutive} if $\{J,J\}\subseteq J$. 
\end{definition}

\AIadd{\formalref{\leandecl{Stafford38/Characteristic/PostScalarExtensionPoisson.lean}{33}{Stafford38.Characteristic.PostScalarExtensionPoisson.IsInvolutive};
negative control \leandecl{Stafford38/Geometry/GeneralCoisotropicSetsTest.lean}{164}{Stafford38.Geometry.GeneralCoisotropicSetsTest.zeroSectionIdealOne_not_isPoisson}}}


\subsection{Gabber's theorem}
\begin{theorem}[Gabber {\cite{Gabber1981}\cite{SinghKumar2014}}{\cite[{Theorem~2.3.1 and Appendix~D}]{HTT}}]\label{ithm:gabber}
Let $R$ be a filtered ring such that $\gr(R)$ is a Noetherian commutative algebra over $\mathbb Q$. Let $M$ be a finitely generated $R$ module. Let $J = \sqrt{Ann_{\gr R}\gr M}$ be the radical of the annihilators of $M$. Then $J$ is involutive, i.e.
\[\{J, J\} \subset J.\]
\end{theorem}


\AIadd{\formalref[{proved in Lean but it is a well known theorem so no need to check the proof}]{%
\leandecl{Stafford38/Characteristic/GabberGlobalAssembly.lean}{224}{Stafford38.Characteristic.GabberGlobalAssembly.weylAssociatedGradedRadicalInvolutivity},
statement \leandecl{Stafford38/Characteristic/CanonicalGabberInvolutivityInterface.lean}{86}{Stafford38.Characteristic.CanonicalGabberInvolutivityInterface.WeylAssociatedGradedRadicalInvolutivity},
minimal primes \leandecl{Stafford38/Characteristic/GabberGlobalAssembly.lean}{144}{Stafford38.Characteristic.GabberGlobalAssembly.minimalPrime_isInvolutive}}}


\section{Mapping to a monic normal form}\label{sec:monic}
Given a fixed nonzero $d \in A_n(k)$ of Bernstein degree $N$, our first goal is to find a map to replace $d$ by a normal form $\tilde{d}$ that is monic in $p$, i.e. of the form $\tilde d=p^N+\sum_{j<N}p^ja_j$ with coefficients $a_j$ in $C$. \\

If $N=0$, then $d \in B_0 = k$ is a nonzero scalar, hence a unit, and $1 = dR+FdS$ holds trivially with $F=1=\ell^0$, $R=d^{-1}$, $S=0$. We therefore assume $N\geq1$.


We want to construct a well-behaved automorphism of $\gr^BA_n(k)=k[x,\xi]$
from simpler linear algebra. We proceed in three stages:
define an automorphism of the vector space $\mathbb V=k^{2n}$, extend it to
an automorphism of $A_n(k)$ itself, and observe the induced action this
extension leaves on the associated graded ring.

We can identify $\mathbb V = k^{2n}$ with the Bernstein-degree-one part $B_1/B_0$ of
$A_n(k)$, i.e.\ linear combinations of $x_1,\ldots,x_n,\partial_1,\ldots,
\partial_n$; the coordinates of a vector are then the coefficients of
these generators. By this identification, $\mathbb V$ carries the standard symplectic form $\omega$ (Def. \ref{def:symplectic-poisson}). 

\begin{lemma}\label{lem:sympcompl}
Let $(\mathbb V,\omega)$ be a symplectic $k$-space of dimension $2n$ and $0\neq v\in \mathbb V$.
Then $v$ is the first vector of a symplectic basis of $\mathbb V$; consequently, there
exists a $k$-linear automorphism $g$ preserving $\omega$ (i.e. $\omega(gu,gv)=\omega(u,v)$ for all
$u,v\in\mathbb V$.) with $g(e)=v$ for any prescribed nonzero vector $e$.
\end{lemma}

\begin{proof}
Since $\omega$ is nondegenerate and $v\neq0$, the functional $\omega(v,-)$
is not identically zero, so there is $w'\in\mathbb V$ with
$\omega(v,w')=c\neq0$; setting $w:=w'/c$ gives $\omega(v,w)=1$.

Let $U=\operatorname{span}(v,w)$. For $0\neq av+bw\in U$: if $a\neq0$,
$\omega(av+bw,w)=a\neq0$; if $b\neq0$, $\omega(av+bw,v)=-b\neq0$. Either
way $av+bw$ pairs nontrivially with something in $U$, so $\omega|_U$ is
nondegenerate.

Since $\omega|_U$ is nondegenerate, $U\cap U^\perp=0$, and as
$\dim U+\dim U^\perp=\dim\mathbb V$ for a nondegenerate form, this gives
$\mathbb V=U\oplus U^\perp$. If $u\in U^\perp$ pairs to zero with all of
$U^\perp$, it also pairs to zero with $U$ (as $u\in U^\perp$), hence with
all of $\mathbb V$, forcing $u=0$; so $\omega|_{U^\perp}$ is nondegenerate
too.

By induction on $n$ (trivial for $n=1$, where $U=\mathbb V$), a
symplectic basis $e_2,\ldots,f_n$ of $U^\perp$ exists; setting $e_1:=v$,
$f_1:=w$ gives a symplectic basis of $\mathbb V=U\oplus U^\perp$, since
cross-pairings between $U$ and $U^\perp$ vanish by definition of
$U^\perp$.

Consequently: given nonzero $v,e\in\mathbb V$, we can complete both vectors to a symplectic basis as above. We then define the linear map $g$ to send $e$ to $v$, and each subsequent $e$-basis vector to the corresponding $v$-basis vector, thus preserving all pairwise values of $\omega$ on basis vectors. Through the bilinearity of $\omega$, $g$ then preserves $\omega$ on all of $\mathbb V$; and $g(e)=v$.
\end{proof}
\AIadd{\formalref{%
\leandecl{Stafford38/Characteristic/PaperSymplecticBasis.lean}{134}{Stafford38.CharacteristicPaperSymplecticBasis.exists_symplectic_basis},
\leandecl{Stafford38/Characteristic/PaperSymplecticBasis.lean}{283}{Stafford38.CharacteristicPaperSymplecticBasis.exists_symplectic_basis_with_pair},
\leandecl{Stafford38/Characteristic/SymplecticCompletion.lean}{238}{Stafford38.CharacteristicSymplecticCompletion.exists_symplectic_mulVec_eq}}}
Let $g$ be a $k$-linear automorphism of $\mathbb V$ that preserves $\omega$.
Since $\mathbb V$ is identified with the span of the generators of $A_n(k)$,
$g$ assigns to each generator $u$ a linear combination $g(u)$ of generators.
Because $A_n(k)$ is presented by generators and the Weyl relations
$[v,u]=\omega(u,v)$, and $g$ preserves $\omega$, the elements $g(u)$ satisfy
the same relations as the $u$. Hence the substitution $u\mapsto g(u)$ extends
to a unique $k$-algebra endomorphism $\Phi_{g}$ of $A_n(k)$. Because $g$ is an automorphism, so is $\Phi_g$, and we have $\Phi_g^{-1} = \Phi_{g^{-1}}.$


\AIadd{\formalref{\leandecl{Stafford38/Weyl/Symplectic.lean}{46}{Stafford38.WeylSymplectic.standardSymplecticAlgHom},
\leandecl{Stafford38/Weyl/Symplectic.lean}{91}{Stafford38.WeylSymplectic.standardSymplecticAlgEquivOfInverse}
(with explicit inverse), \leandecl{Stafford38/Weyl/SymbolCompatibility.lean}{416}{Stafford38.WeylSymbolCompatibility.standardSymplecticAlgHom_principal_compatibility}}}

\begin{proposition}\label{prop:monic}
Let $0\neq d\in A_n(k)$ with $N:=\deg_B d\ge1$. There exists an automorphism $\Phi_g$ of $A_n(k)$ and
$c\in k^{\times}$ such that
\begin{equation}
  \tilde d:=c^{-1}\Phi_g(d)=p^{\,N}+\sum_{j<N}p^{\,j}a_j, 
  \qquad a_j\in C,\quad \deg_B a_j\le N-j\label{eq:monic-d}
\end{equation}
in the notation of Convention~\ref{conv:pair}. $\tilde d$ itself has $\deg_B \tilde d = N.$
\end{proposition}

\begin{proof}
Let $s\in k[x,\xi]$ be the Bernstein principal symbol of $d$, a nonzero
homogeneous polynomial of degree $N$ on $\mathbb V$. A characteristic-zero field is
infinite, so every nonzero polynomial has a nonvanishing value: we choose $v\in \mathbb V$
with $s(v)\neq0$. Because $s$ is homogeneous of degree $N\geq 1$, $s(0)$ must vanish, therefore $v$ is nonzero. 
Let $e_p\in\mathbb V$ be the standard basis vector corresponding to $p$.
By Lemma~\ref{lem:sympcompl} there is an $\omega$-preserving automorphism
$h$ of $\mathbb V$ with $h(e_p)=v$. The standard basis is symplectic, so
the Gram matrix $J$ of $\omega$ in it satisfies $J^{-1}=-J$; hence
$h^{\mathsf T}Jh=J$ implies $hJh^{\mathsf T}=J$, i.e.\ the transpose
$g:=h^{\mathsf T}$ is again $\omega$-preserving. Let $\Phi_g$ be the
corresponding automorphism of $A_n(k)$, extending $u\mapsto g(u)$ on
generators.
Let $M$ be the matrix whose rows are the images of the generators, i.e.\
$g(u_i)=\sum_j M_{ij}u_j$. Since $\Phi_g$ sends each generator to a
degree-one element, it maps $B_m$ into $B_m$ for all $m$, and since
Bernstein symbols are multiplicative, a standard monomial of degree $N$ is
mapped to an element of $B_N$ whose class modulo $B_{N-1}$ is obtained by
substituting $u_i\mapsto\sum_j M_{ij}u_j$ in the corresponding commutative
monomial. Hence $\Phi_g(d)\equiv s(Mu)\bmod B_{N-1}$. As $M$ is invertible,
$s(Mu)$ is again a nonzero homogeneous polynomial of degree $N$, so
$\deg_B\Phi_g(d)=N$ and the Bernstein symbol of $\Phi_g(d)$ is the
function $s_g\colon a\mapsto s(Ma)$. Because $g=h^{\mathsf T}$,
the rows of $M$ are the columns of $h$, i.e.\ $Me_j=h(e_j)$. Hence the
value of $s_g$ at $e_p$ is
\[ s_g(e_p) = s(Me_p)=s(h(e_p))=s(v)\neq0 .\]
Set $c:=s(v)$. Since $c$ is a nonzero scalar, it is a central unit of
$A_n(k)$, and $\tilde d:=c^{-1}\Phi_g(d)$ satisfies $\deg_B\tilde d=N$ and
$\sigma^B(\tilde d)(e_p)=1$.
Using the normal form $A_n(k)=\bigoplus_j p^{\,j}C$, we can write
$\tilde d=\sum_{j = 0}^N p^{\,j}a_j$ with $a_j\in C$. The Bernstein filtration is
compatible with this decomposition, so
$N=\deg_B\tilde d=\max_j\,(j+\deg_B a_j)$, and we have $\deg_B a_j\le N-j$ for
all $j$; in particular $a_j=0$ for $j>N$ and $a_N\in C\cap B_0=k$. 
The symbol is
\[
  \sigma^B(\tilde d)=\sum_{j+\deg_B a_j=N}\sigma^B(p)^{\,j}\,\sigma^B(a_j),
\]
For $j<N$ the factors $\sigma^B(a_j)$ must have a positive degree and must therefore be non-constant in a coordinate other than $p$, so they vanish at $e_p$. Because we defined $\tilde d$ to be $c^{-1}\Phi_g(d)$ , evaluating at $e_p$ gives $a_N=\sigma^B(\tilde d)(e_p)=c^{-1}\sigma^B(\Phi_g(d))(e_p) = c^{-1}s_g(e_p) = c^{-1}s(v)=1$, and therefore
\[ \tilde d=p^{\,N}+\sum_{j<N}p^{\,j}a_j,\qquad a_j\in C,\quad
  \deg_B a_j\le N-j. \]
\end{proof}

\AIadd{\formalref[{Actual form is not proved in lean, what is important and subsequently used is the fact that the symbol is monic in p and its coefficient  is 1. ("PBWmonicity")}]{%
\leandecl{Stafford38/Weyl/MonicNormalization.lean}{158}{Stafford38.WeylMonicNormalization.exists_normalized_symplectic_image},
\leandecl{Stafford38/Weyl/MonicNormalization.lean}{181}{Stafford38.WeylMonicNormalization.scalar_or_normalized_symplectic_image},
\leandecl{Stafford38/Weyl/PBWMonicBridge.lean}{98}{Stafford38.WeylPBWMonicBridge.IsPBWMonicAt},
\leandecl{Stafford38/Weyl/PaperRightMonic.lean}{669}{Stafford38.WeylPaperRightMonic.presentedOuter_right_monic_decomposition},
\leandecl{Stafford38/Characteristic/SymplecticCompletion.lean}{346}{Stafford38.CharacteristicSymplecticCompletion.exists_symplectic_chart_matrices}}}


\begin{remark}\label{rem:transport}
Since $\Phi_g$ is a ring automorphism, so an identity $1=\tilde dR+x_1^N\tilde dS$ would transport back to
\[ 1=d\bigl(c^{-1}\Phi_g^{-1}(R)\bigr)+F\,d\bigl(c^{-1}\Phi_g^{-1}(S)\bigr),
  \qquad F=\Phi_g^{-1}(x_1)^N.\]
\end{remark}

\AIadd{\formalref[{compare}]{\leandecl{Stafford38/UniversalAssembly.lean}{69}{Stafford38.UniversalAssembly.universalStatement_of_canonicalSupportVanishing},
\leandecl{Stafford38/FixedSourceAssembly.lean}{24}{Stafford38.FixedSource.universalFixedSourceStatement_of_canonicalSupportVanishing}}}

We now use the order filtration on $A_n(k)$ (where coordinates have
degree zero and derivatives have degree one) to characterize $P =
\sigma_{\mathrm{ord}}(\tilde d)$, the order principal symbol of $\tilde
d$. Recall that a symbol, being the class of an element in a single
graded piece, is automatically \emph{homogeneous} of the corresponding
degree, i.e.\ contains only terms of that total degree.
\begin{definition}\label{def:monic-xi1}
We say a polynomial $P=\sum_{i=0}^Nc_i\,\xi_1^i\in
k[\xi_1,\ldots,\xi_n]$, written as a polynomial in $\xi_1$ with
coefficients $c_i \in k[\xi_2,\ldots,\xi_n]$, is \emph{monic of degree $N$ in
$\xi_1$} if $c_N=1$.
\end{definition}
\AIadd{\formalref{\texttt{Polynomial.Monic} after \leandecl{Stafford38/Characteristic/NormalSymbolPolynomial.lean}{28}{Stafford38.Characteristic.NormalSymbolPolynomial.normalSymbolAlgEquiv}}}

\begin{lemma}\label{lem:order-symbol}
Let $\tilde d=p^N+\sum_{j<N}p^{\,j}a_j$ with $a_j\in C$, $\deg_Ba_j\le N-j$,
as in \eqref{eq:monic-d}, and set $N=\deg_B\tilde d$. Then $\tilde d$ has
order $N$, and its order principal symbol is a polynomial in $\xi$ alone,
\begin{equation} \sigma_{\mathrm{ord}}(\tilde d)=P(\xi),\qquad
 P\in k[\xi_1,\ldots,\xi_n],\qquad \deg P=N.
 \label{eq:P}\end{equation}
$P$ is monic of degree $N$ in $\xi_1$ and in particular
\begin{equation} P(1,0,\ldots,0)=1.
 \label{eq:P-e1}\end{equation}
\end{lemma}

\begin{proof}
Since $\deg_B\tilde d=N$, every PBW monomial $x^\alpha\partial^\beta$
occurring in $\tilde d$ satisfies $|\alpha|+|\beta|\le N$. A monomial
contributing to the order-$N$ part has $|\beta|=N$, hence $|\alpha|=0$: it
is a pure $\partial^\beta$ with scalar coefficient. Thus $\tilde d$ has
order $N$, and its order-$N$ part is a $k$-linear combination of
$\partial^\beta$ with $|\beta|=N$, giving $\sigma_{\mathrm{ord}}(\tilde
d)=P(\xi)$, homogeneous of degree $N$ with coefficients in $k$.

The monomial $p^N=\partial_1^N$ occurs in $\tilde d$ with coefficient $1$;
every other order-$N$ monomial comes from some $p^ja_j$ with $j<N$ and
$a_j\in C$ (hence involving no further factor of $p=\partial_1$), so it
carries a positive power of some transverse $\partial_i$, $i>1$. Such a
monomial contributes only to coefficients $c_i$ with $i<N$, never to
$c_N$; hence $c_N$ receives no contribution besides the coefficient of
$p^N$ itself, giving $c_N=1$ identically, i.e.\ $P$ is monic of degree $N$
in $\xi_1$. Evaluating at $(1,0,\ldots,0)$ then gives $P(1,0,\ldots,0)=1$.
\end{proof}
\AIadd{\formalref{\leandecl{Stafford38/Characteristic/InitialIdeal.lean}{320}{Stafford38.CharacteristicInitialIdeal.canonical_orderPrincipalComponent_isFibreOnly},
\leandecl{Stafford38/Geometry/ConormalAxisContradiction.lean}{61}{Stafford38.Geometry.ConormalAxisContradiction.canonical_orderPrincipalComponent_isHomogeneous},
\leandecl{Stafford38/Characteristic/NormalSymbolPolynomial.lean}{104}{Stafford38.Characteristic.NormalSymbolPolynomial.canonicalNormalPolynomial_monic},
\leandecl{Stafford38/Geometry/ConormalAxisContradiction.lean}{131}{Stafford38.Geometry.ConormalAxisContradiction.canonical_orderPrincipalComponent_eval_pureMomentumAxis}}}


\section{$x_1$ as a surjective map}\label{sec:euler}
We define $E=x_1p$ and
\begin{equation} I=\tilde dA_n+x_1^N\tilde dA_n,\qquad Q=A_n/I.
 \label{eq:I-Q}\end{equation}
The goal of this section will be to show that $Q = Qx_1$ (Proposition \ref{prop:x-surjective}).
\AIadd{\formalref{\leandecl{Stafford38/Weyl/EulerResidue.lean}{26}{Stafford38.WeylEulerResidue.canonicalRightIdeal}
$=\operatorname{span}\{d,x^Nd\}$ as a right ideal}}

The commutator of $E$ defines a grading of $A_n$: 
$x_1$ has Euler degree $1$, $p$ has Euler degree $-1$, and the transverse Weyl algebra
has Euler degree zero. Let $A_{\geq0}$ denote the nonnegative part of $A_n$ in this grading. In PBW form (with $x_i$ on the left of $\partial_i$), $A_{\geq0}$ is generated by the transverse algebra, $E$, and $x_1$. 

%\AIadd{Indeed $A_n$ has the basis $b\,x_1^ap^c$ with $b$ transverse, of Euler degree $a-c$; if $a\ge c$ then $x_1^ap^c=x_1^{a-c}(x_1^cp^c)$ and $x_1^cp^c$ is a polynomial in $E$ by Lemma~\ref{lem:euler}. In particular every element of Euler degree $\ge1$ lies in $A_{\ge0}x_1$, using $Ex_1=x_1(E+1)$ to move $x_1$ to the right.}
By the commutation relations of the Weyl algebra, it holds that
\begin{align}
    Ex_1=x_1(E+1),\qquad x_1E=(E-1)x_1,\nonumber \\
    Ep=p(E-1), \qquad pE=(E+1)p
    \label{eq:nonnegative-pbw}
\end{align}
Additionally, $x_1$ commutes with the transverse algebra, so $x_1$ is normal in $A_{\geq 0}.$ This also means that every element of Euler degree $\ge1$ lies in $A_{\ge0}x_1$. 
\AIadd{\formalref[{Lean uses the subring generated by the transverse algebra, $x_1$ and $x_1p$, and one-sided normality $x_1r=r'x_1$, which is all that is used}]{%
\leandecl{Stafford38/Weyl/EulerSubring.lean}{25}{Stafford38.WeylEulerSubring.eulerSubring},
\leandecl{Stafford38/Weyl/EulerSubring.lean}{71}{Stafford38.WeylEulerSubring.coordinate_normal_eulerSubring}}}

\begin{lemma}\label{lem:euler}
For every $m\ge0$,
\begin{equation}\label{eq:eulerprod}
  x_1^{m}p^{\,m}=\prod_{i=0}^{m-1}(E-i), \qquad p^{\,m} x_1^{m}=\prod_{i=1}^{m}(E+i),
\end{equation}
the factors in each product commuting with one another.
\end{lemma}

\begin{proof}
Since both products in \eqref{eq:eulerprod} are polynomials in the single element $E$, their factors commute.
We induct on $m$; both identities are trivial for $m=0$. Assuming the
first identity for $m$,
\[\begin{aligned}
x_1^{m+1}p^{\,m+1}
 &=x_1\Bigl(\prod_{i=0}^{m-1}(E-i)\Bigr)p
 =\Bigl(\prod_{i=0}^{m-1}(E-1-i)\Bigr)x_1p\\
 &=\Bigl(\prod_{i=1}^{m}(E-i)\Bigr)E
 =\prod_{i=0}^{m}(E-i),
\end{aligned}\]
using $x_1E=(E-1)x_1$ repeatedly and the definition of $E$. The second identity follows
identically from $Ep=p(E-1)$.
%\[p^{m+1}x_1^{m+1}=p\Bigl(\prod_{i=1}^m(E+i)\Bigr)x_1=\Bigl(\prod_{i=1}^m(E+1+i)\Bigr)px_1=\Bigl(\prod_{i=2}^{m+1}(E+i)\Bigr)(E+1).\]
\end{proof}
\AIadd{\formalref{\leandecl{Stafford38/Weyl/EulerProductIdentities.lean}{70}{Stafford38.Weyl.EulerProductIdentities.eval_fallingEulerProduct} and
\leandecl{Stafford38/Weyl/EulerProductIdentities.lean}{79}{Stafford38.Weyl.EulerProductIdentities.eval_risingEulerProduct} ,
\leandecl{proofs/weyl_pure_power.lean}{148}{Stafford.exists_eulerPolynomial_x_pow_mul_d_pow}}}

\begin{lemma}\label{lem:euler-residue}
There exists an element $U\in A_{\geq0}$ such that
\begin{equation} 1+Ux_1\in I.
 \label{eq:euler-residue}\end{equation}
\end{lemma}

\begin{proof}
From Lemma \ref{lem:euler} we define
\[p^Nx_1^N=\prod_{r=1}^N(E+r)=:R_+(E),\qquad
 x_1^Np^N=\prod_{r=0}^{N-1}(E-r)=:R_-(E).\]
The roots of $R_+$ and $R_-$ are disjoint (the sets $\{-1,\ldots,-N\}$ and $\{0,\ldots,N-1\}$ of integers, embedded in the prime field
$\mathbb Q\subset k$), so $\gcd(R_+,R_-)=1$ in $\mathbb Q[E]$. By the
extended Euclidean algorithm, there exist $u,v\in\mathbb Q[E]\subset
k[E]$ with
\[ R_+(E)u(E)+R_-(E)v(E)=p^Nx_1^Nu(E)+x_1^Np^Nv(E)=1.\]
We write $\tilde d=p^N+\Sigma$ with $\Sigma = \sum_{j<N}p^{\,j}a_j$. Then
\[(\tilde d\,x_1^N)u(E)+(x_1^N\tilde d)v(E)=1+w\in I,\]
where
\[w=\Sigma x_1^Nu(E)+x_1^N\Sigma v(E).\]
We can write each $a_j$ in $\Sigma$ in powers of $x_1$, with coefficients in the transverse Weyl algebra. A term $p^jx_1^rb$ coming from $p^ja_j$ has $r\geq0$ and $j<N$, and multiplication by $x_1^N$ gives that term Euler degree at least $N-j>0$. $u$ and $v$ only depend on $E$, so they have Euler degree 0. Therefore every term of $w$ has positive Euler degree, and thus $w\in A_{\geq0}x_1$ - or equivalently, there exists a $U\in A_{\geq0}$, such that $w=Ux_1$.
\end{proof}
\AIadd{\formalref[{the ring identity is Lean's; the factorization uses left coefficients $c\,x^ap^j$ but that does not change anything because Euler degree does not depend on the order of x and p}]{%
\leandecl{Stafford38/Weyl/EulerResidue.lean}{41}{Stafford38.WeylEulerResidue.positiveEulerResidue_identity},
\leandecl{Stafford38/Weyl/EulerRemainder.lean}{257}{Stafford38.WeylEulerRemainder.positiveEulerResidue_eq_eulerSubring_mul_coordinate},
\leandecl{Stafford38/Weyl/EulerRemainder.lean}{405}{Stafford38.WeylEulerRemainder.presentedPositiveEulerResidue}}}


\begin{proposition}\label{prop:x-surjective}
Right multiplication by $x_1$ on $Q$ is surjective.
\end{proposition}
\begin{proof}
Let $q$ be the class of $1$ in $Q$. By Equation \eqref{eq:euler-residue}, we have
$q=-qUx_1$, so $q\in(qA_{\geq0})x_1$. We use the normality of $x_1$ to show
\begin{equation} qA_{\geq0}=(qA_{\geq0})x_1.
 \label{eq:normal-surj}\end{equation}
The inclusion $\supseteq$ holds because $x_1\in A^{\geq0}$ and
$A^{\geq0}$ is a ring. For $\subseteq$, let $r\in A^{\geq0}$ and choose
$r'$ with $x_1r=r'x_1$ (possible since $x_1$ is normal) then $qr=-(qU)x_1r=-(qUr')x_1\in(qA^{\geq0})x_1$ since $Ur'\in A^{\geq0}$.

$qA_{\geq0}$ is stable under right multiplication by $p$: If $m=m'x_1$, then
\[ mp=m'(x_1p)=m'E\in qA_{\geq0}.\]
It is of course also stable under $C$, since $C=A_{n-1}(k)[x_1]\subseteq A_{\ge0}$ because no element of $C$ contains $p$. Since $C$ and $p$ together generate $A_n$, we have $qA_{\geq0}=qA_n=Q$. Equation \eqref{eq:normal-surj} is therefore $Q=Qx_1$.
\end{proof}
\AIadd[\formalref{The proof in Lean is element wise and not over the closedness of Ageq0, but it is so elementary that i think it is okay to leave this one as it is. Annahmen und outcome sollte das gleiche sein. }]{\formalref{\leandecl{Stafford38/Weyl/QuotientTransport.lean}{240}{Stafford38.WeylQuotientTransport.presentedCanonicalRightQuotient_rightMul_coordinate_surjective},
\leandecl{Stafford38/Quotient/EulerSurjectivity.lean}{86}{Stafford38.EulerSurjectivity.rightMul_surjective_of_euler_normal};
generation by $C$ and $p$: \leandecl{Stafford38/Weyl/EulerRemainder.lean}{310}{Stafford38.WeylEulerRemainder.closure_pairOldSubring_union_coordinate_momentum_eq_top}}}


\section{Vanishing cokernel}\label{sec:cokernel}
\begin{definition}[{\cite{Stacks}[Tag 00CM]}] \label{def:localization}
Let $R$ be a commutative ring and $S\subseteq R$ a multiplicative subset
(i.e.\ $1\in S$ and $S$ is closed under products). The \emph{localization}
$S^{-1}R$ is the ring of fractions $a/s$ ($a\in R$, $s\in S$), where
\[ \frac{a}{s}=\frac{a'}{s'}\iff t(s'a-sa')=0\ \text{ for some }t\in S,\]
with the usual addition and multiplication of fractions. For an $R$-module
$N$, the \emph{localization} $S^{-1}N$ is the $S^{-1}R$-module of fractions
$y/s$ ($y\in N$, $s\in S$), with $y/s=y'/s'$ iff $t(s'y-sy')=0$ for some
$t\in S$. For a prime $\mathfrak p\subset R$ we write
$R_{\mathfrak p}=S^{-1}R$ and $N_{\mathfrak p}=S^{-1}N$ for
$S=R\setminus\mathfrak p$. Localization is exact: it sends exact sequences
of $R$-modules to exact sequences of $S^{-1}R$-modules.
\end{definition}
\begin{definition}[{\cite{Stacks}[Definition 10.40.1]}]\label{def:supp}
    Let $R$ be a ring and let $M$ be an $R$-module. The \emph{support} of $M$ is the set 
    \[\Supp_R(M) = \{\mathfrak{p} \in \Spec(R) | M_{\mathfrak{p}} \neq 0\}.\]
\end{definition}

For this section, we write $M=\operatorname{gr}Q$ and denote the right multiplication by $x_1$ by $f:M\rightarrow M$. We furthermore define
\[
 \begin{gathered}
 V=\ker(f),\qquad U=\operatorname{coker}(f) =M/x_1M\\
 T=k[x_2,\ldots,x_n,\xi_2,\ldots,\xi_n],
 \end{gathered}
\]
hence $T[x_1,\xi_1] = k[x,\xi]$. If $n=1$, we define $T=k$.
We want to prove that $U=0$ by comparing
two lengths over a localization of $T$.
\begin{definition}[\mbox{\cite{coutinho, cohn}, \cite[Tags~02LY, 00J3, 00IV]{Stacks}}]\label{def:length}
    The length $\ell_T$ of a $T$-module of finite length $N$ is defined as the number of steps in a composition series of $N$ such that
    \[0 = N_0\subsetneq N_1\subsetneq \dots \subsetneq N_{\ell_T} = N\]
    where each successive quotient $N_i/N_{i-1}$ is simple (i.e. has no proper nontrivial submodules). The length is well-defined, as it does not depend on the choice of the composition series. Lengths are additive for short exact sequences.
\end{definition}
\AIadd{\formalref{\texttt{Module.length} (Mathlib); finite length of the localized modules:
\leanlib{AlgebraicAnalysis/Module/MinimalSupportKernelCokernelLengths.lean}{48}{AlgebraicAnalysis.localized_kernel_and_cokernel_isFiniteLength}}}
\begin{lemma}\label{lem:supp-criterion}
Let $R$ be a commutative ring, $N$ an $R$-module and $\mathfrak p\subset R$ a
prime. Then $\mathfrak p\in\operatorname{Supp}_R(N)$ if and only if there is
$y\in N$ with $\operatorname{Ann}_R(y)\subseteq\mathfrak p$.
\end{lemma}
\AIadd{[\formalref{Lemma used in subsequent proof, possibly not explicitly mentioned in Lean}]}
\begin{proof}
In $N_{\mathfrak p}$, a
fraction is zero when it can be cleared by an element outside $\mathfrak p$.
For $y/1$ this means: there is $u\in R\setminus\mathfrak p$ with $uy=0$.
Such a $u$ lies in $\operatorname{Ann}_R(y)$ and not in $\mathfrak p$. So
\[y/1=0\iff\operatorname{Ann}_R(y)\not\subseteq\mathfrak p,\]
and, taking the contrapositive,
\[y/1\neq0\iff\operatorname{Ann}_R(y)\subseteq\mathfrak p.\]
\emph{$(\Leftarrow)$} Suppose $y\in N$ has
$\operatorname{Ann}_R(y)\subseteq\mathfrak p$. Then, $y/1\neq0$ in
$N_{\mathfrak p}$. So $N_{\mathfrak p}$ has a nonzero element, hence
$N_{\mathfrak p}\neq0$ and $\mathfrak p\in\operatorname{Supp}_R(N)$.\\
\emph{$(\Rightarrow)$} Suppose $N_{\mathfrak p}\neq0$. Pick a nonzero
element of $N_{\mathfrak p}$. It has the form $y/s$ with $y\in N$ and
$s\notin\mathfrak p$. The element $s$ is invertible in $N_{\mathfrak p}$, so
\[y/s=0\iff y/1=0.\]
Since $y/s\neq0$, we get $y/1\neq0$, and therefore
$\operatorname{Ann}_R(y)\subseteq\mathfrak p$, which is the desired $y$.
\end{proof}
\begin{lemma}\label{lem:tangential-finiteness}
The modules $U$ and $V$ are finite over $T$, and
$\operatorname{Supp}_T V\subseteq\operatorname{Supp}_T U$.
If $\mathfrak q$ is a minimal prime in $\operatorname{Supp}_T U$, then
$U_{\mathfrak q}$ and $V_{\mathfrak q}$ have finite length over
$T_{\mathfrak q}$.
\end{lemma}

\begin{proof}
The good quotient filtration makes $M = \operatorname{gr} Q$ finitely generated over $\operatorname{gr} A_n = k[x,\xi] = T[{x_1}, \xi_1]$.
Since $P$ annihilates $M$ (Corollary \ref{cor:charP}), this action factors
through $T[{x_1}, \xi_1]/(P)$, which, because $P$ is monic of degree $N$ in $\xi_1$, is generated over $T[{x_1}]$ by
$1,\xi_1,\ldots,\xi_1^{N-1}$. Thus $M$ is finite over $T[{x_1}]$.
Because $M$ is finite, and $T[{x_1}]$ is a Noetherian ring, the kernel $V$ and cokernel $U$ of multiplication in $M$ by ${x_1}$ are also finite over $T[{x_1}]$. Since they are annihilated by ${x_1}$, they are finite over $T$. %Noetherian theorem 3.5 https://www.cambridge.org/core/services/aop-cambridge-core/content/view/5DD13688B6CC1628C3FD7A318CD479FE/9781139171762c1_p1-19_CBO.pdf/commutative-rings-and-modules.pdf

Let $\mathfrak p\in\operatorname{Supp}_T(V)$. By Lemma \ref{lem:supp-criterion}, we can choose $v\in V=\ker f$ with
$\operatorname{Ann}_T(v)\subseteq\mathfrak p$. Let us now consider the annihilators of $v$ in $k[x,\xi]=T[x_1, \xi_1]$, defining $S$ as the image of $T\setminus \mathfrak p$ in $k[x,\xi]$: $\operatorname{Ann}_{k[x,\xi]}(v)\cap S = \varnothing$, because were $r \in T\setminus\mathfrak p$ and in $\operatorname{Ann}_{k[x,\xi]}(v)$, then $v$ would be annihilated by $r$, which in turn would imply that $r$ is also in $\operatorname{Ann}_T(v)$. Since $\operatorname{Ann}_T(v) \subseteq \mathfrak p$, but we assumed $r \in T\setminus\mathfrak p$, this is a contradiction. There must therefore exist a prime $\mathfrak q \in k[x,\xi]$ with $\mathfrak q \cap S = \varnothing$ and $\operatorname{Ann}_{k[x,\xi]}(v) \subseteq \mathfrak q$. Since $v \in V = \ker f$, by Lemma \ref{lem:supp-criterion}, it follows that $\mathfrak q \in \Supp_{k[x,\xi]}(\ker f).$
From Vasconcelos' theorem \cite[Theorem~2.4]{matsumura} follows that $\operatorname{Supp}_{k[x,\xi]}(\ker f)\subseteq \operatorname{Supp}_{k[x,\xi]}(\operatorname{coker}f)$: 
We prove the contrapositive. Let $\mathfrak q$ be a prime of ${k[x,\xi]}$ with
$\mathfrak q\notin\operatorname{Supp}_{k[x,\xi]}(\operatorname{coker}f)$, i.e.\
$(\operatorname{coker}f)_{\mathfrak q}=0$.
Since localization is exact, $(\operatorname{coker}f)_{\mathfrak q}=
\operatorname{coker}(f_{\mathfrak q})$, so
$\operatorname{coker}(f_{\mathfrak q})=0$: the endomorphism
$f_{\mathfrak q}\colon M_{\mathfrak q}\to M_{\mathfrak q}$ is therefore surjective.
The localization $M_{\mathfrak q}$ is a finite module over the Noetherian ring
${k[x,\xi]}_{\mathfrak q}$, and $f_{\mathfrak q}$ is ${k[x,\xi]}_{\mathfrak q}$-linear, because
$f$ is multiplication by $x_1\in {k[x,\xi]}$ and ${k[x,\xi]}$ is commutative. By Vasconcelos'
theorem, a surjective endomorphism of a finite module is injective. Hence
$f_{\mathfrak q}$ is injective, and by exactness of localization,
$(\ker f)_{\mathfrak q}=\ker(f_{\mathfrak q})=0$, i.e.\
$\mathfrak q\notin\operatorname{Supp}_{k[x,\xi]}(\ker f)$. The contrapositive is therefore $\mathfrak q\in\operatorname{Supp}_{k[x,\xi]}(\ker f)$,
then $\mathfrak q\in\operatorname{Supp}_{k[x,\xi]}(\operatorname{coker}f) = \Supp_{k[x,\xi]}(U).$ By
Lemma~\ref{lem:supp-criterion} there therefore exists an $u\in U$ with
$\operatorname{Ann}_{k[x,\xi]}(u)\subseteq\mathfrak q$. 
To get back to the support over $T$, we define $r\in\operatorname{Ann}_T(u)$. Then $ru=0$, so
$r\in\operatorname{Ann}_{k[x,\xi]}(u)\subseteq\mathfrak q$. If $r\notin\mathfrak p$,
the image of $r$ would lie in $\mathfrak q\cap S=\varnothing$, which is
impossible. Hence $r\in\mathfrak p$, i.e.\
$\operatorname{Ann}_T(u)\subseteq\mathfrak p$, and
Lemma~\ref{lem:supp-criterion} gives
$\mathfrak p\in\operatorname{Supp}_T(U)$. Since $\mathfrak p\in
\operatorname{Supp}_T(V)$ was arbitrary,
$\operatorname{Supp}_T(V)\subseteq\operatorname{Supp}_T(U)$.
\end{proof}
\AIadd{\formalref{\leandecl{Stafford38/Characteristic/CanonicalOldTangentialFiniteness.lean}{73}{Stafford38.Characteristic.CanonicalOldTangentialFiniteness.canonical_finite_old_coordinate_kernel_cokernel},
\leandecl{Stafford38/Characteristic/CanonicalTangentialSymbolFiniteness.lean}{116}{Stafford38.Characteristic.CanonicalTangentialSymbolFiniteness.canonical_orderAssociatedGradedModule_finite_tangential_coordinate}}}

\AIcomment{ALG-04}{\textbf{Max --- correspondence check; no rewrite of this lemma requested.} Check the monic quotient argument giving finite $T$-modules $U,V$, together with the subsequent $T$-action and $T$-linear page maps. Localization applies to the associated graded pages. See COMP-06.}
\begin{lemma}\label{lem:page-inequality}
For every minimal prime $\mathfrak q$ in $\operatorname{Supp}_T U$,
\[
 \ell_{T_{\mathfrak q}}(U_{\mathfrak q})
 \le \ell_{T_{\mathfrak q}}(V_{\mathfrak q}).
\]
\end{lemma}

\begin{proof}
To prove the inequality we set up a spectral sequence measuring how $\ker(f)$ and
$\operatorname{coker}(f)$ acting on $Q$ itself compare, order by order, to $V$ and
$U$, acting on $\gr Q$.

Consider the filtered complex $Q\xrightarrow{f}Q$, with decreasing filtration $G^p=F_{-p}Q$, where $F$ is
the order filtration. Since ${x_1}$ has order $0$, $f(G^p)\subseteq G^p$ for every
$p$; the induced map on $G^p/G^{p+1}$, summed over $p$, is exactly the action of
${x_1}$ on $M=\operatorname{gr}Q$. This is the crudest approximation to $\ker_Q(f)$ and
$\operatorname{coker}_Q(f)$: it correctly identifies $V,U$ as approximations, but it
does not yet see elements that vanish to higher order than this naive count
predicts. The spectral sequence refines this approximation in successive stages,
called \emph{pages}: for each $r\ge0$, the source and target pages are
\[
 A_r=\bigoplus_{p\in\mathbb Z}\frac{G^p\cap f^{-1}(G^{p+r})}{G^{p+1}\cap f^{-1}(G^{p+r})},
 \qquad
 B_r=\bigoplus_{p\in\mathbb Z}\frac{G^p}{(G^p\cap f(G^{p-r+1}))+G^{p+1}},
\]
where the degree-$p$ summand of $A_r$ isolates elements of $G^p$ whose image under
$f$ drops $r$ extra filtration levels (candidates for the kernel, accurate to
order $r$), and the degree-$p$ summand of $B_r$ isolates the part of $G^p$ not yet
hit by $f$ from $r-1$ levels shallower.

Applying $f$ to a representative gives a well-defined map $d_r:A_r\to B_r$,
shifting the summand index by $r$; one checks directly from these definitions that
\[ A_0=B_0=M,\qquad d_0={x_1},\qquad A_{r+1}=\ker d_r,\qquad B_{r+1}=\operatorname{coker}d_r,\]
so that each page is obtained from the previous one by taking kernel and cokernel
of $d_r$, exactly as in the construction of a spectral sequence from a filtered
complex.

These pages carry a $T$-module structure. The pages $A_r,B_r$ are
subquotients of $Q$, which a priori carries an action of the transverse Weyl
algebra $A_{n-1}(k)$ by right multiplication. This action commutes with
$f$, since the transverse algebra commutes with $x_1$, so the maps
$d_r$ are compatible with it. Two transverse elements commute only up to a
commutator, which lowers the order filtration by one. On the page
quotients such lower-order terms vanish, so the induced actions commute, and
the commutative ring $T$ acts on every $A_r$ and $B_r$, via the symbols of
the transverse generators. The maps $d_r$ and the successor-page maps are
$T$-linear. In particular, this means that localizing at $T\setminus\mathfrak q$ and
measuring lengths over $T_{\mathfrak q}$ is a well-defined action for the pages.

Now we fix a minimal prime $\mathfrak q$ in $\operatorname{Supp}_TU$ and localize at
$T\setminus\mathfrak q$. We localize the $T$-modules
$A_r,B_r$ and their $T$-linear maps. Exactness of localization preserves the
successor-page kernel and cokernel identities. By
Lemma~\ref{lem:tangential-finiteness}, $(A_1)_{\mathfrak q}=V_{\mathfrak q}$ and
$(B_1)_{\mathfrak q}=U_{\mathfrak q}$ have finite length over $T_{\mathfrak q}$; since each
$(A_{r+1})_{\mathfrak q}\subseteq (A_r)_{\mathfrak q}$ and $(B_{r+1})_{\mathfrak q}$ is a quotient of $(B_r)_{\mathfrak q}$, all later pages
have finite length as well.


The maps $(B_1)_{\mathfrak q}\to (B_r)_{\mathfrak q}$ are surjective by construction.
Let $K_r:=\ker((B_1)_{\mathfrak q}\to (B_r)_{\mathfrak q})$. Since $(B_1)_{\mathfrak q}\to (B_{r+1})_{\mathfrak q}$ factors through $(B_1)_{\mathfrak q}\to (B_r)_{\mathfrak q}\to (B_{r+1})_{\mathfrak q}$, we have
$K_r\subseteq K_{r+1}$, so $(K_r)_{r\ge1}$ is an increasing chain of submodules of $(B_1)_{\mathfrak q}$.
This chain exhausts $(B_1)_{\mathfrak q}$: Before localization, take a homogeneous
class in $B_1$, represented by $z\in G^p$. Since $Q=Qx_1$, choose one preimage
$z=wx_1$. Exhaustiveness of the filtration means that $w$ must be in $G^s$ for some $s$; choose $r\ge1$ with
$p-r+1\le s$. Then $w\in G^{p-r+1}$, so $z$ belongs to the denominator of
$B_r$ and its class vanishes. Every class is a finite sum of homogeneous
classes and hence vanishes on a sufficiently late page. The same holds for a
localized class $z/a$, since its numerator already vanishes on such a page.
As $z\in B_1$ was arbitrary, so is $z/a$ in $(B_1)_{\mathfrak q}$, and $\bigcup_{r\ge1}K_r=(B_1)_{\mathfrak q}$.
Since $(B_1)_{\mathfrak q}$ has finite length, it is a Noetherian $T_{\mathfrak q}$-module. By the ascending
chain condition, the increasing chain $K_1\subseteq K_2\subseteq\cdots$ stabilizes: there exists a finite $R$ with
$K_R=K_{R+i}$ for all $i\geq 0$. Combined with $\bigcup_rK_r=(B_1)_{\mathfrak q}$, stabilization forces $K_R=(B_1)_{\mathfrak q}$.
Finally, $K_R=(B_1)_{\mathfrak q}$ means $(B_1)_{\mathfrak q}\to (B_R)_{\mathfrak q}$ must be the zero map: by the first isomorphism theorem,
$(B_1)_{\mathfrak q}/K_R\cong\operatorname{im}((B_1)_{\mathfrak q}\to (B_R)_{\mathfrak q})$, and $K_R=(B_1)_{\mathfrak q}$ gives $(B_1)_{\mathfrak q}/K_R=0$, so
$\operatorname{im}((B_1)_{\mathfrak q}\to (B_R)_{\mathfrak q})=0$. But $(B_1)_{\mathfrak q}\to (B_R)_{\mathfrak q}$ is also surjective, so $(B_R)_{\mathfrak q}=0$.\\

At each stage $r\ge1$ the differential $(d_r)_{\mathfrak q}$ gives an exact sequence
\[ 0\to (A_{r+1})_{\mathfrak q}\to (A_{r})_{\mathfrak q}\xrightarrow{(d_r)_{\mathfrak q}}(B_{r})_{\mathfrak q}\to (B_{r+1})_{\mathfrak q}\to0,\]
which splits at $\operatorname{im}((d_r)_{\mathfrak q})$ into two short exact sequences.
Additivity of length gives
$\ell((A_{r})_{\mathfrak q})=\ell((A_{r+1})_{\mathfrak q})+\ell(\operatorname{im}(d_r)_{\mathfrak q})$ and
$\ell((B_{r})_{\mathfrak q})=\ell(\operatorname{im}(d_r)_{\mathfrak q})+\ell((B_{r+1})_{\mathfrak q})$; eliminating
$\ell(\operatorname{im}(d_r)_{\mathfrak q})$ yields, for every $r\ge1$,
\begin{equation}\label{eq:additivity}
 \ell((B_r)_{\mathfrak q})+\ell((A_{r+1})_{\mathfrak q})=\ell((A_{r})_{\mathfrak q})+\ell((B_{r+1})_{\mathfrak q}).
\end{equation}

Summing this identity over $r=1,\ldots,R-1$ and reindexing the sums, most terms cancel on both sides, leaving % $\sum_{r=1}^{R-1}\ell((A_{r+1})_{\mathfrak q})=\sum_{r=2}^{R}\ell((A_{r})_{\mathfrak q})$,
%$\sum_{r=1}^{R-1}\ell((B_{r+1})_{\mathfrak q})=\sum_{r=2}^{R}\ell((B_{r})_{\mathfrak q})$, the terms $\ell((A_{r})_{\mathfrak q}),\ell((B_{r})_{\mathfrak q})$ for $2\le r\le R-1$
%cancel on both sides, leaving
\[\ell((B_1)_{\mathfrak q})+\ell((A_R)_{\mathfrak q})=\ell((A_{1})_{\mathfrak q})+\ell((B_{R})_{\mathfrak q}).\]
Since $(B_{R})_{\mathfrak q}=0$ and $(A_{1})_{\mathfrak q}=V_{\mathfrak q}$, $(B_{1})_{\mathfrak q}=U_{\mathfrak q}$, this reads
\[ \ell(U_{\mathfrak q})+\ell((A_{R})_{\mathfrak q})=\ell(V_{\mathfrak q}).\]
As $\ell((A_{R})_{\mathfrak q})\ge0$, we conclude
\[ \ell_{T_{\mathfrak q}}(U_{\mathfrak q})\le\ell_{T_{\mathfrak q}}(V_{\mathfrak q}).\]
\end{proof}

\AIadd{[{\href{https://itpplasma.github.io/stafford38-formal/lean_proof_details.pdf\#nameddest=comparison-COMP-06}{Readable correspondence note.}}]\formalref{\leandecl{Stafford38/Characteristic/CanonicalPageEulerInequality.lean}{25}{Stafford38.Characteristic.CanonicalPageEulerInequality.canonicalPage_length_target_le_source};
exhaustion \leanlib{AlgebraicAnalysis/Module/UniformBoundaryVanishing.lean}{106}{AlgebraicAnalysis.exists_uniform_subsingleton_localized}}}


\noindent On the other hand, we can also establish a different inequality.

\begin{lemma}\label{lem:strict-support-inequality}
For every minimal prime $\mathfrak q$ in $\operatorname{Supp}_T U$,
\[
 \ell_{T_{\mathfrak q}}(U_{\mathfrak q})
 > \ell_{T_{\mathfrak q}}(V_{\mathfrak q}).
\]
\end{lemma}
\begin{proof}
Let us consider $\operatorname{Ann}_C M$. By Gabber's Theorem (\ref{ithm:gabber}, \cite{Gabber1981}), $I = \sqrt{\operatorname{Ann}_C M}$ is involutive. 
If $\mathfrak p$ is a minimal prime of $J$,
then $\mathfrak p$ is itself involutive, i.e. $\{\mathfrak p,\mathfrak p\}\subseteq\mathfrak p$. We can show this the following way:
We localize $k[x,\xi]$ at $\mathfrak p$. Since $\mathfrak p$ is minimal over $I$, the only prime of $k[x,\xi]_{\mathfrak p}$ containing $I_{\mathfrak p}$ is $\mathfrak p_{\mathfrak p}$ itself; as $I$ is already radical, so is its localization, giving
\[I_{\mathfrak p} = \sqrt{I_{\mathfrak p}} = \mathfrak p_{\mathfrak p}.\]
The Poisson bracket extends to $k[x,\xi]_{\mathfrak p}$ since by the Leibniz rule, $\{a/s,\,b/t\}$ is determined by $\{a,b\}$ and the brackets of $a,b$ with $s,t$, giving a well-defined bracket on $k[x,\xi]_\mathfrak p$ agreeing with the one on $k[x,\xi]$ under $f\mapsto f/1$; in particular $\{f/1,g/1\}=\{f,g\}/1$ for $f,g\in k[x,\xi]$, and the involutivity of $I$ localizes to $\{I_\mathfrak p,I_\mathfrak p\}\subseteq I_\mathfrak p$.
Now take $f,g\in\mathfrak p$. Then $f/1,g/1\in\mathfrak p_\mathfrak p=I_\mathfrak p$, so
\[ \{f,g\}/1 = \{f/1,g/1\} \in I_\mathfrak p = \mathfrak p_\mathfrak p.\]
Thus there is $t\notin\mathfrak p$ with $t\{f,g\}\in\mathfrak p$. Since $\mathfrak p$ is prime and $t\notin\mathfrak p$, this forces $\{f,g\}\in\mathfrak p$.

\AIadd{\formalref{%
\leandecl{Stafford38/Characteristic/PaperLocalizedInvolutivity.lean}{187}{Stafford38.Characteristic.PaperLocalizedInvolutivity.minimalPrimes_isInvolutive_of_radical_isInvolutive_viaLocalization}}}

From the involutivity of $\mathfrak p$ follows that no $\mathfrak p$ contains ${x_1}$: Assume ${x_1}\in \mathfrak p$. Since $P \in \mathfrak p$,  also the Poisson brackets
\[ \{{x_1}, P\}, \quad \{{x_1}, \{x_1, P\}\}, \dots\]
lie in $\mathfrak p$. $\{{x_1},-\}$ differentiates once with respect to $\xi_1$, and because $P$ is homogeneous of degree $N$ and monic in $\xi_1$, its $N$th derivative is the nonzero constant $N!$ - therefore the $N$th bracket would put $N!$ in $\mathfrak p$. Since $\mathfrak p$ is a proper prime ideal, no nonzero constant can be in $\mathfrak p$, so we must conclude ${x_1} \notin \mathfrak p$.
This result holds under localization at $T\setminus\mathfrak q$: no minimal prime of $\operatorname{Ann}_{k[x,\xi]_\mathfrak p}(M_{\mathfrak q})$ contains ${x_1}$
either, since the annihilator of the finite module $M$ commutes with localization.
Consider the increasing chain of kernels of powers of the localization $f_{\mathfrak q}$ of right multiplication with $x_1$
$\operatorname{ker}(f_{\mathfrak q})$ $\subseteq \operatorname{ker}(f^2_{\mathfrak q})$ $\subseteq \dots$.
Since $M_{\mathfrak q}$ is Noetherian, this chain stabilizes. Let
$K:=\ker(f^m_{\mathfrak q})$ at a stage where it stabilizes, so $f_{\mathfrak q}(K) = {x_1}K\subseteq K$ and
$\ker(f^{m+1}_{\mathfrak q})=\ker(f_{\mathfrak q}^m)=K$. Each quotient $\ker(f_{\mathfrak q}^{i+1})/\ker(f_{\mathfrak q}^i)$ maps to
$V_{\mathfrak q}=\ker(f_{\mathfrak q})$ injectively via $v\mapsto {x_1}^iv$, so, as $V_{\mathfrak q}$ has
finite length, so does each quotient and hence $K$ itself.

Two elementary facts about $K$:
\begin{itemize}
\item ${f_{\mathfrak q}}^{-1}K=K$: if $v\in K$ then $f_{\mathfrak q}(v)\in K$ (as $f_{\mathfrak q}(K)\subseteq K$); if $f_{\mathfrak q}(v)\in K=\ker(f_{\mathfrak q}^m)$
then $f_{\mathfrak q}^{m+1}(v)=0$, so $v\in\ker({f_{\mathfrak q}}^{m+1})=K$.
\item $K\cap {x_1}M_{\mathfrak q}={x_1}K = f_{\mathfrak q}(K)$: "$\supseteq$" is trivial. For "$\subseteq$", if $w=f_{\mathfrak q}(v)\in K$ for
some $v\in M_{\mathfrak q}$, then $v\in f_{\mathfrak q}^{-1}K=K$, so $w\in f_{\mathfrak q}(K) = x_1K$.
\end{itemize}

For any endomorphism $\varphi$ of a finite-length module $N$,
additivity along $0\to\ker \varphi\to N\to \varphi(N)\to0$ and $0\to \varphi(N)\to N\to N/\varphi(N)\to0$ gives
$\ell(\ker \varphi)=\ell(N/\varphi(N))$. We can apply this to $\varphi=f_{\mathfrak q}|_K$: since $\ker(f_{\mathfrak q})\subseteq K$, we
have $\ker(f_{\mathfrak q}|_K)=\ker(f_{\mathfrak q})=V_{\mathfrak q}$, so with $f_{\mathfrak q}(K) = x_1K,$
\[\ell(K/{x_1}K)=\ell(V_{\mathfrak q}).\]

On the other hand, the natural surjection $M_{\mathfrak q}/({x_1}M_{\mathfrak q})\to (M_{\mathfrak q}/K)/({x_1}(M_{\mathfrak q}/K))$ has
the kernel $(K+({x_1}M_{\mathfrak q}))/({x_1}M_{\mathfrak q})\cong K/(K\cap ({x_1}M_{\mathfrak q}))=K/({x_1}K)$, giving (we now write $E = (M_{\mathfrak q}/K)$ for better readability) the exact sequence
\[0\longrightarrow K/({x_1}K)\longrightarrow U_{\mathfrak q}\longrightarrow
 E/({x_1}E)\longrightarrow0,\]
so by additivity and using $\ell(K/({x_1}K))=\ell(V_{\mathfrak q})$,
\[\ell(U_{\mathfrak q}) = \ell(V_{\mathfrak q}) + \ell\bigl(E/({x_1}E)\bigr).\]
It remains to show the last term is positive.

Since $\mathfrak q\in\operatorname{Supp}_TU$,
$U_{\mathfrak q}=M_{\mathfrak q}/({x_1}M_{\mathfrak q})\ne0$. We choose $\mathfrak P\in\operatorname{Supp}_{k[x,\xi]_\mathfrak p}(M_{\mathfrak q}/({x_1}M_{\mathfrak q}))$
and a minimal prime $\mathfrak p\subseteq\mathfrak P$ of $M_{\mathfrak q}$. We have shown above that ${x_1}\notin\mathfrak p$, so ${x_1}$ is a unit in the localization at $\mathfrak p$; as $K$ is annihilated by a power of ${x_1}$, $K_{\mathfrak p}=0$. From
$0\to K_{\mathfrak p}\to (M_{\mathfrak q})_{\mathfrak p}\to E_{\mathfrak p}\to0$ we get
$E_{\mathfrak p}\cong (M_{\mathfrak q})_{\mathfrak p}\ne0$, so $\mathfrak p$ lies in the support of $M_{\mathfrak q}/K$ in $k[x,\xi]_\mathfrak p$.
Since support is closed under passing to larger primes \cite[Lemma 10.40.5 (Tag 00L2)]{Stacks} and
$\mathfrak p\subseteq\mathfrak P$, also $\mathfrak P\in\operatorname{Supp}_{k[x,\xi]_\mathfrak p}E$, i.e.
$E_{\mathfrak P}\ne0$.

Since $\mathfrak P\in\operatorname{Supp}_{k[x,\xi]_\mathfrak p}(M_{\mathfrak q}/({x_1}M_{\mathfrak q}))$ and $M_{\mathfrak q}/({x_1}M_{\mathfrak q})$ is annihilated by
${x_1}$, we have ${x_1}\in\mathfrak P$. If $({x_1}E_{\mathfrak P})=E_{\mathfrak P}$, then as
${x_1}$ lies in the maximal ideal of the local ring $(k[x,\xi]_\mathfrak p)_{\mathfrak P}$, Nakayama's lemma \cite{nakayama}, (see also \cite[Tag~00DV]{Stacks})
would force $E_{\mathfrak P}=0$, a contradiction. Hence
$\bigl(E/({x_1}E)\bigr)_{\mathfrak P}\ne0$, so $E/({x_1}E)\ne0$ and $\ell\left(E/({x_1}E)\right) > 0$.

Therefore $\ell(U_{\mathfrak q})=\ell(V_{\mathfrak q})+\ell\bigl(E/{x_1}E\bigr)
>\ell(V_{\mathfrak q})$.
\end{proof}
\AIadd{\formalref{\leanlib{AlgebraicAnalysis/Module/BaseLocalizedKoszulPositivity.lean}{25}{AlgebraicAnalysis.BaseLocalizedKoszulPositivity.localized_length_cokernel_gt_kernel};
$x_1$ in no minimal prime: \leandecl{Stafford38/Characteristic/NoncharacteristicMinimalPrime.lean}{118}{Stafford38.Characteristic.NoncharacteristicMinimalPrime.canonical_minimalPrime_mem_of_normalCoordinate_false}}}


\begin{lemma}\label{lem:Uzero}
The cokernel $U=M/x_1M$ of right multiplication by $x_1$ on
$M=\operatorname{gr}Q$ vanishes, i.e. $U=0$.
\end{lemma}
\begin{proof}
If $U\ne0$, it has a minimal support prime $\mathfrak q$.
Lemma~\ref{lem:tangential-finiteness} makes the relevant lengths finite,
and Lemmas~\ref{lem:page-inequality} and
\ref{lem:strict-support-inequality} contradict one another.
Thus $U = M/{x_1}M=0$.
\end{proof}
\AIadd{\formalref{\leandecl{Stafford38/Characteristic/CanonicalKoszulContradiction.lean}{86}{Stafford38.Characteristic.CanonicalKoszulContradiction.coordinate_cokernel_subsingleton},
\leandecl{Stafford38/Characteristic/CanonicalKoszulContradiction.lean}{50}{Stafford38.Characteristic.CanonicalKoszulContradiction.no_minimal_coordinate_cokernel_support}}}


\section{Characteristics of the characteristic variety}\label{sec:char}
\begin{corollary}\label{cor:charP}
Equip $Q$ with the order filtration $Q_m:=qF_m=\operatorname{im}(F_m\to Q)$
induced by the cyclic generator $q = [1]$, the class of $1$ in $Q$. As before, we denote the symbol of $\tilde d$ in the order filtration by  $P=\sigma_N(\tilde d)$. Then
\begin{equation}\label{eq:P-annihilates}
 P\in\Ann_{k[x,\xi]}\gr Q.
\end{equation}
\end{corollary}

\begin{proof}
$Q_\bullet$ is a good filtration because $Q$ is cyclic. Let $\bar q$ be the class of $q$ in $\gr Q$. We know that $q\,\tilde d=[\tilde d]=0$ in $Q$ (because $\tilde d \in I$) and $\tilde d$ has order $N$ (Lemma \ref{lem:order-symbol}), so $q\tilde d \in Q_N$. Therefore $\bar q\cdot P$ , its image in $Q_N/Q_{N-1}$, must also be zero, and hence $P$ annihilates the class $\bar q$ of $q$ in $\gr Q$.
Since $\gr Q$ is generated over $k[x,\xi]$ by $\bar q$, we get $P\in\Ann_{k[x,\xi]}\gr Q$. 
\end{proof}

\AIadd{\formalref{\leandecl{Stafford38/Characteristic/InitialIdeal.lean}{281}{Stafford38.CharacteristicInitialIdeal.canonical_orderPrincipalComponent_mem_initialIdeal},
\leandecl{Stafford38/Characteristic/InitialIdeal.lean}{291}{Stafford38.CharacteristicInitialIdeal.canonical_orderCharacteristicSupport_subset_principal_zeroLocus}}}
\begin{proposition}\label{prop:char-avoids}
$\Char(Q)\cap\pi^{-1}(H)=\varnothing.$
\end{proposition}
\begin{proof}
Write $M=\gr Q$. Since $M$ is finite over $k[x,\xi]$,
$\Supp(M/x_1M)=\Supp(M)\cap V(x_1)$ \cite[Tag~00L3]{Stacks}, where $V(x_1)$ is
the set of primes of $k[x,\xi]$ containing $x_1$. By Lemma~\ref{lem:Uzero},
$M/x_1M=0$, so its support is empty, hence $\Supp(M)\cap V(x_1)=\varnothing$.
By definition $\Supp(M)=\Char(Q)$, and $V(x_1)=\pi^{-1}(H)$ is the preimage of
$H=\{x_1=0\}$ under the projection $\pi$. Therefore
$\Char(Q)\cap\pi^{-1}(H)=\varnothing$.
\end{proof}

\begin{proof}
We write $M = \gr Q$. Since $M$ is finite over ${k[x,\xi]}$, its quotient by ${x_1}$ has support $\Supp_{k[x,\xi]}(M/x_1M) = \operatorname{Supp}_{k[x,\xi]}(M)\cap V({x_1})$ \cite[Tag~00L3]{Stacks}. By Lemma \ref{lem:Uzero} $M/{x_1}M=0$ and has therefore no support, so $\Supp_{k[x,\xi]}(M)\cap V(x_1) = \varnothing$. Per definition $\operatorname{Supp}_{k[x,\xi]}(M) = \Char(Q)$, and $V(x_1)=\pi^{-1}(H)$ is the preimage of $H=\{x_1=0\}$ under the projection $\pi$. Therefore
$\Char(Q)\cap\pi^{-1}(H)=\varnothing$.

\end{proof}

\AIadd{\formalref[{stated for the support transported along $\xi\mapsto-\xi$, which fixes $x_1$; transported back verbatim}]{%
\leandecl{Stafford38/Characteristic/CanonicalKoszulContradiction.lean}{123}{Stafford38.Characteristic.CanonicalKoszulContradiction.canonical_support_avoidance},
\leandecl{Stafford38/Characteristic/SpecializedNoncharacteristicEquality.lean}{151}{Stafford38.SpecializedNoncharacteristicEquality.transposedSupport_disjoint_axis_iff}}}





\section{A conormal direction at infinity}\label{sec:conormal-review}
\begin{comment}
\subsection{Geometric vocabulary for this section}\label{sec:geom-vocab}

Throughout this section $K$ denotes an algebraically closed field of
characteristic zero.  We write
$X=\mathbb A^n_K$, use the notation of Section~\ref{sec:cotangent}, and
``point'' means closed point.

We recall the standard notions we need (see \cite[Ch.~I]{hartshorne}), and
then fix the notions specific to this section in numbered definitions.

A subset $Y\subseteq X$ is \emph{(Zariski) closed} if it is the common zero
set $V(J)$ of a set of polynomials, and \emph{irreducible} if it is not the
union of two proper closed subsets. Its \emph{vanishing ideal}
$I(Y)=\{f\in K[x]: f|_Y=0\}$ is radical, and prime if and only if $Y$ is
irreducible; by the Nullstellensatz, $I(V(J))=\sqrt J$. For irreducible $Y$,
$R=K[x]/I(Y)$ is its coordinate ring, $F=\operatorname{Frac}(R)$ its
function field, and $\dim Y=\operatorname{trdeg}_KF$. The \emph{Zariski
closure} $\overline S$ of a subset $S$ is the smallest closed set containing
$S$; a polynomial vanishing on $S$ vanishes on $\overline S$.

The tangent space $T_yX=K^n$ is the space of directions at a point
$y\in X$, and a covector $\xi=\sum_i\xi_i\,dx_i\in T_y^*X=(K^n)^*$ acts on
$v\in K^n$ by $\xi(v)=\sum_i\xi_iv_i$. For a polynomial $f$ and a point $y$,
the \emph{differential} of $f$ at $y$ is the covector
\[
 df_y=\sum_{i=1}^n\frac{\partial f}{\partial x_i}(y)\,dx_i\in T_y^*X,
\]
so that $df_y(v)$ is the derivative of $f$ at $y$ in the direction $v$; it
is linear in $f$, satisfies $d(fg)_y=f(y)\,dg_y+g(y)\,df_y$, and gives
$d(x_i)_y=dx_i$.

Projective space $\mathbb P^{n-1}=\mathbb P((K^n)^*)$ is the set of nonzero
covectors modulo nonzero scalars; $[\xi]$ denotes the class of $\xi$, so that
$[dx_1]=[1:0:\dots:0]$. A homogeneous polynomial $P(\xi)$ has a well-defined
zero set $V_+(P)\subseteq\mathbb P((K^n)^*)$, which is closed. We write
$[Z_0:\dots:Z_n]$ for homogeneous coordinates on $\mathbb P^n$.


We write $\kappa[[t]]$ and $\kappa((t))$ for formal power series and formal
Laurent series; a power series is a \emph{unit} if and only if its constant
term is nonzero, and \emph{evaluating} (or \emph{reducing}) at $t=0$ means
taking the constant term. A \emph{formal arc} is a map
$\operatorname{Spec}K[[t]]\to Z$; its \emph{generic point} lies over
$K((t))$ ($t\neq0$) and its \emph{special point} over $t=0$. A closed
subvariety $Z$ of a projective variety contains the special point of an arc
as soon as it contains its generic point (valuative criterion of properness,
\cite[Thm.~II.4.7]{hartshorne}). A vector over $\kappa[[t]]$ is
\emph{primitive} if its entries are not all divisible by $t$. A matrix over
$\kappa[[t]]$ whose reduction at $t=0$ has full column rank has a left
inverse over $\kappa[[t]]$; its columns then span a \emph{direct summand},
and the annihilator $L^\perp$ of a direct summand $L$ reduces onto
$(L\bmod t)^\perp$.

A \emph{$K$-derivation}
$D\colon\kappa\to\kappa$ is a $K$-linear map with $D(ab)=aD(b)+bD(a)$, and
$\Omega_{\kappa/K}$ is the $\kappa$-vector space of differentials, dual to the
derivations. In characteristic zero every field extension is separable, so
$\dim_\kappa\Omega_{\kappa/K}=\operatorname{trdeg}_K\kappa$, and the
differentials of a transcendence basis form a basis. Derivations act on
power series coefficientwise.

\begin{definition}[Smooth points, tangent and conormal spaces]\label{def:smooth-conormal}
Let $Y\subseteq X$ be closed and irreducible of dimension $m$. A point
$y\in Y$ is \emph{smooth} if the differentials $df_y$, $f\in I(Y)$, span a
space of dimension $n-m$; the smooth points form a dense
open subset $Y_{\mathrm{sm}}\subseteq Y$. At $y\in Y_{\mathrm{sm}}$ the
\emph{tangent space} and the \emph{conormal space} are
\[
 T_yY=\bigcap_{f\in I(Y)}\ker df_y\subseteq K^n,\qquad
 N_y^*Y=\operatorname{span}\{df_y: f\in I(Y)\}\subseteq(K^n)^*;
\]
$N_y^*Y$ is the space of covectors vanishing on $T_yY$, of dimension $n-m$.
\end{definition}

\begin{definition}[Conormal cone and conormal variety]\label{def:conormal}
For $Y$ as above, the \emph{conormal cone} is the union of all conormal
spaces, forgetting the base points,
\[
 N(Y)=\bigcup_{y\in Y_{\mathrm{sm}}}N_y^*Y\subseteq(K^n)^*;
\]
it is closed under scaling, but in general not closed. The \emph{conormal
variety} $T^*_YX\subseteq T^*X$ is the Zariski closure of
$\bigcup_{y\in Y_{\mathrm{sm}}}\{y\}\times N_y^*Y$; it keeps the base points
and, being closed, also contains the limiting conormal covectors at the
singular points of $Y$.
\end{definition}

\begin{definition}[Projective closure]\label{def:proj-closure}
We embed $X\hookrightarrow\mathbb P^n$ as the \emph{affine chart}
$\{Z_0\neq0\}$ via $x_i=Z_i/Z_0$; its complement $\{Z_0=0\}$ is the
\emph{hyperplane at infinity}. The \emph{projective closure}
$\overline Y\subseteq\mathbb P^n$ of $Y$ is the Zariski closure of $Y$, so
$\overline Y\smallsetminus Y\subseteq\{Z_0=0\}$ (as $Y$ is closed in $X$).
On a chart $\{Z_j\neq0\}$ we use the ratios $q_i=Z_i/Z_j$. The \emph{affine
cone} over a projective linear subspace is the corresponding linear subspace
of $K^{n+1}$, and the \emph{projective tangent plane} at a smooth point of
$\overline Y$ is the projective linear subspace whose affine cone is spanned
by the position vector $q$ and the partial derivatives of a local
parametrization.
\end{definition}



\begin{definition}[Divisorial valuation ring]\label{def:divisorial}
Let $Y$ be irreducible with function field $F$. A DVR $V$ with
$K\subseteq V\subset F$ and $\operatorname{Frac}(V)=F$ is \emph{divisorial} if
it is the local ring of a prime divisor on some normal variety with function
field $F$; then $\operatorname{trdeg}_K\kappa=\dim Y-1$ for its residue field
$\kappa$.
\end{definition}
\begin{definition}
An element of a field extension is \emph{algebraic}
over a subfield if it satisfies a nonzero polynomial equation with
coefficients in that subfield, and that an extension is algebraic if all its
elements are.
\end{definition}


\AIadd{\formalref[{$Y$ given by a prime ideal; conormal space = annihilator of the Zariski tangent space; closure by homogeneous polynomials}]{%
\leandecl{Stafford38/Geometry/ProjectiveConormalDirections.lean}{40}{Stafford38.Geometry.ProjectiveConormalDirections.smoothConormalDirectionSet},
\leandecl{Stafford38/Geometry/ProjectiveConormalDirections.lean}{32}{Stafford38.Geometry.ProjectiveConormalDirections.projectiveHomogeneousClosure}}}
\end{comment}
\begin{comment}

\AIadd{\formalref[{generic point only Zariski-tangent; no smoothness needed for this step}]{%
\leandecl{Stafford38/Geometry/ConormalScalarExtensionVanishing.lean}{137}{Stafford38.Geometry.ConormalScalarExtensionVanishing.scalarExtension_vanishing};
\leandecl{Stafford38/Geometry/SmoothConormalFibreVanishing.lean}{24}{Stafford38.Geometry.SmoothConormalFibreVanishing.fibreLift_mem_vanishingIdeal_equationConormal};
\leandecl{Stafford38/Geometry/LaurentConormalDirection.lean}{71}{Stafford38.Geometry.LaurentConormalDirection.residue_mem_extensionFibreClosure_of_laurent_generic};
\leandecl{Stafford38/Geometry/GeneralTangentLimitCriterion.lean}{501}{Stafford38.Geometry.GeneralTangentLimitCriterion.exists_axis_laurent_smooth_conormal_direction}}}
\end{comment}
\begin{comment}
Throughout this section $K$ is an algebraically closed field of characteristic
zero, $X=\mathbb A^n_K$ with $T^*X=X\times(K^n)^*$ as in
Section~\ref{sec:cotangent}, and $H=\{x_1=0\}$. We identify a covector with its
coordinate vector, so that $df_y=\nabla f(y)=(\partial f/\partial x_1(y),\dots,
\partial f/\partial x_n(y))$ for a polynomial $f$. We write $[Z_0:\dots:Z_n]$ for
homogeneous coordinates on $\mathbb P^n$ and embed $X\hookrightarrow\mathbb P^n$
as the chart $\{Z_0\neq0\}$ via $x_i=Z_i/Z_0$. Let $Y\subseteq X$ be a nonempty
closed irreducible subvariety with vanishing ideal $I(Y)$. A point $y\in Y$ is
\emph{smooth} if the $df_y$, $f\in I(Y)$, span a space of dimension
$n-\dim Y$; the smooth points form a dense open subset $Y_{\mathrm{sm}}$. For
$y\in Y_{\mathrm{sm}}$ the \emph{conormal space} is
$N_y^*Y=\operatorname{span}\{df_y:f\in I(Y)\}\subseteq(K^n)^*$, the \emph{conormal
cone} is $N(Y)=\bigcup_{y\in Y_{\mathrm{sm}}}N_y^*Y$, and $\operatorname{Dir}(Y)$ is
the Zariski closure in $\mathbb P((K^n)^*)$ of $\{[\xi]:0\neq\xi\in N(Y)\}$.

\begin{theorem}\label{thm:asymptotic-conormal-formal}
Let $Y\subseteq X$ be nonempty, closed and irreducible with $Y\cap H=\varnothing$.
\begin{enumerate}[label=(\roman*)]
\item Every polynomial $P\in K[\xi_1,\dots,\xi_n]$ vanishing on $N(Y)$ satisfies
$P(1,0,\dots,0)=0$.
\item $[dx_1]\in\operatorname{Dir}(Y)$.
\end{enumerate}
\end{theorem}

Let $\mathfrak p=I(Y)$, which is prime, so $Y=V(\mathfrak p)$. We write
$R=K[x]/\mathfrak p$, $F=\operatorname{Frac}(R)$, and $\bar x_i$ for the image of
$x_i$ in $R$. Since $Y\cap H=\varnothing$, the Nullstellensatz
(Theorem~\ref{thm:nss}) gives $(\mathfrak p,x_1)=K[x]$, so $\bar x_1\in R^\times$.

If $\bar x_1$ is algebraic over $K$, there is a nonzero polynomial $g$ over $K$
in one variable with $g(\bar x_1)=0$, i.e.\ $g(x_1)\in\mathfrak p$. Since $K$ is
algebraically closed, $g(x_1)=\lambda\prod_i(x_1-c_i)$ with $\lambda\in K^\times$,
and as $\mathfrak p$ is prime, some factor $x_1-c$ lies in $\mathfrak p$. Thus
$x_1=c$ on $Y$, and $c\neq0$ because $Y\ne\varnothing$ avoids $x_1=0$. In this
case $f=x_1-c\in\mathfrak p=I(Y)$ has differential $df_y=dx_1$ at every
$y\in Y_{\mathrm{sm}}$, so $dx_1\in N_y^*Y$ and its coordinate vector
$(1,0,\dots,0)$ lies in $N(Y)$. Hence every polynomial vanishing on $N(Y)$
vanishes at $(1,0,\dots,0)$, and $[dx_1]\in\operatorname{Dir}(Y)$ (see the end of
the proof). So we assume from now on that $\bar x_1$ is transcendental over $K$.

The rest of the proof follows this rough logic. We build an arc to a place at
infinity where $x_1$ vanishes, with generic point $y$ on $Y$
(Lemma~\ref{lem:place}). We construct a basis of the tangent cone of $Y$ at $y$
that depends on $t$ and stays a basis at $t=0$ (Lemma~\ref{lem:tangent}), and show
that its limit lies in $\{Z_1=0\}$. Hence the covector $Z_1^\vee$ reading off the
$Z_1$-coordinate annihilates the limit, and we lift it to a family of covectors
$\ell(t)$ annihilating the tangent cones for $t\neq0$, with $\ell(0)=Z_1^\vee$
(Lemma~\ref{lem:lift}). For $t\neq0$ the $\ell(t)$ are conormal covectors to $Y$
at $y$, so their limit, whose affine part is $dx_1$, lies in the closure of the
conormal directions (Lemma~\ref{lem:special}). We first prove the lemmas, then
the theorem.

\medskip\noindent
For a discrete valuation ring $V\subset F$ with maximal ideal
$\mathfrak m_V=(\varpi)$, every nonzero $f\in F$ is $f=\varpi^ku$ with $k\in\mathbb Z$
and $u\in V^\times$, and we write $\operatorname{ord}_V(f)=k$.

\begin{lemma}\label{lem:place}
Let $\bar x_1$ be transcendental over $K$ and put $Z_0=1$, $Z_i=\bar x_i$
$(1\le i\le n)$. There are a discrete valuation ring $V$ with
$K\subseteq V\subset F$, $\operatorname{Frac}V=F$, and a column
$q=(q_0,\dots,q_n)$ with $q_i=sZ_i$ for some $s\in F^\times$, such that, with
$\kappa=V/\mathfrak m_V$, $\bar q_i$ the residue of $q_i$, $a=\operatorname{ord}_V(q_0)$
and $b=\operatorname{ord}_V(q_1)$:
\begin{enumerate}[label=(\alph*)]
\item $q_i\in V$ for all $i$, $q_j=1$ for some $j$, and $q_1=q_0\bar x_1$;
\item $\bar x_1\in\mathfrak m_V$, $R\not\subseteq V$, and $0<a<b$; in particular
$j\notin\{0,1\}$;
\item $\kappa$ is algebraic over $L=K(\bar q_0,\dots,\bar q_n)$;
\item $d:=\dim_\kappa\Omega_{\kappa/K}$ is finite, and there are indices
$i_1,\dots,i_d\notin\{0,1,j\}$ such that $d\bar q_{i_1},\dots,d\bar q_{i_d}$ form a
basis of $\Omega_{\kappa/K}$; let $D_1,\dots,D_d$ be the dual $K$-derivations of
$\kappa$, so $D_k(\bar q_{i_l})=\delta_{kl}$;
\item $\Omega_{V/K}$ is a finitely generated $V$-module.
\end{enumerate}
\end{lemma}
\begin{proof}
\emph{Step 1: a valuation ring.} Since $\bar x_1$ is transcendental over $K$,
$S=K[\bar x_1]$ is a polynomial ring and $\bar x_1S$ is a proper ideal. By
Chevalley's extension theorem there is a valuation ring $W\subseteq F$ with
$S\subseteq W$ and $\bar x_1\in\mathfrak m_W$; in particular $K\subseteq W$ and
$\bar x_1^{-1}\notin W$.

\emph{Step 2: the column.} For nonzero $a',b'\in F$, either $a'/b'\in W$ or
$b'/a'\in W$. So among the finitely many nonzero $Z_a$ we can choose $j$ with
$Z_a/Z_j\in W$ for all $a$. Put $s=Z_j^{-1}$ and $q_i=sZ_i$. Then $q_j=1$,
$q_i\in W$, and $q_1/q_0=Z_1/Z_0=\bar x_1$. As $\bar x_i=q_i/q_0$, the ring
$C_0=K[q_0,\dots,q_n,\bar x_1]\subseteq W$ is finitely generated with
$\operatorname{Frac}C_0=F$.

\emph{Step 3: the divisor.} Let $A$ be the integral closure of $C_0$ in $F$; it is
a finite $C_0$-module (finiteness of normalization), hence a finitely generated
integrally closed Noetherian domain with $\operatorname{Frac}A=F$ and
$q_i\in A$. As $W$ is integrally closed, $A\subseteq W$, so $\bar x_1^{-1}\notin A$
and $\bar x_1$ is a nonzero nonunit of $A$. By Krull's principal ideal theorem
there is a prime $\mathfrak p_D\ni\bar x_1$ of $A$ of height one, and
$V:=A_{\mathfrak p_D}$ is a discrete valuation ring with
$K\subseteq A\subseteq V$, $\operatorname{Frac}V=F$ and $\bar x_1\in\mathfrak m_V$.
This gives (a) and (e): $\Omega_{A/K}$ is a finite $A$-module and
$\Omega_{V/K}=V\otimes_A\Omega_{A/K}$.

\emph{Step 4: (b).} If $R\subseteq V$, then $\bar x_1\in R^\times\subseteq V^\times$,
contradicting $\bar x_1\in\mathfrak m_V$. Next, $b-a=\operatorname{ord}_V(q_1/q_0)=
\operatorname{ord}_V(\bar x_1)>0$, and $a,b\ge0$ as $q_0,q_1\in V$. If $a=0$, then
$q_0\in V^\times$, so all $\bar x_i=q_i/q_0\in V$ and $R\subseteq V$, a
contradiction. Thus $0<a<b$, and $j\notin\{0,1\}$ since $q_j=1$ is a unit while
$q_0,q_1$ are not.

\emph{Step 5: (c).} $\kappa=\operatorname{Frac}(A/\mathfrak p_D)$. As $A$ is integral
over $C_0$ and $\bar x_1\in\mathfrak p_D$, the ring $A/\mathfrak p_D$ is integral over
the image $K[\bar q_0,\dots,\bar q_n]$ of $C_0$, so $\kappa$ is algebraic over $L$
(algebraic elements form a field).

\emph{Step 6: (d).} $\kappa$ is finitely generated over $K$, so $d<\infty$. Since
$\bar q_0=\bar q_1=0$ (as $a,b>0$) and $\bar q_j=1$, the field $L$ is generated over
$K$ by the $\bar q_i$ with $i\notin\{0,1,j\}$. Choose among them a transcendence
basis $B=\{\bar q_{i_1},\dots,\bar q_{i_d}\}$ of $L/K$; by (c) it is one of
$\kappa/K$. As $\kappa/K(B)$ is algebraic, hence separable in characteristic zero,
$d\bar q_{i_1},\dots,d\bar q_{i_d}$ is a basis of $\Omega_{\kappa/K}$ (see e.g.\
\cite[II, \S8]{hartshorne}). The dual basis of
$\operatorname{Hom}_\kappa(\Omega_{\kappa/K},\kappa)=\operatorname{Der}_K(\kappa,\kappa)$
gives the $D_k$.
\end{proof}

\noindent\textbf{Setup.} From now on fix $V,q,j,\kappa,a,b,d,i_l,D_k$ as in
Lemma~\ref{lem:place}, and let $\varpi$ be a generator of $\mathfrak m_V$. Let
$\widehat V$ be the completion of $V$. Since $V\supseteq K$ and $K$ has
characteristic zero, there exists a coefficient field of $\widehat V$
containing $K$ and mapping isomorphically onto $\kappa$ \cite[Tag~0C0S]{Stacks};
choosing one gives an isomorphism $\widehat V\cong\kappa[[t]]$ under which
$\varpi$ corresponds to the formal variable $t$. The natural map
$V\to\widehat V$ is injective ($\bigcap_k\varpi^kV=0$) and extends to an embedding
$F\hookrightarrow\kappa((t))$, through which we identify elements of $F$ with
Laurent series. The constant term of a series is its residue, so
$q_i(t)=\sum_rc_{ir}t^r$ with $c_{i0}=\bar q_i$, and in $\kappa[[t]]$
\[q_0=t^au_0,\qquad q_1=t^bu_1,\]
with $u_0(0)\neq0\neq u_1(0)$. Dividing by $q_0\neq0$, we define
$y=(q_1/q_0,\dots,q_n/q_0)\in\kappa((t))^n$; its coordinates are the images of
$\bar x_1,\dots,\bar x_n$ under $F\hookrightarrow\kappa((t))$. Since ring
homomorphisms commute with polynomial evaluation and this map is injective, for
every polynomial $f\in K[x_1,\dots,x_n]$ we have
\begin{equation}\label{eq:generic}
 f(y)=0\iff f(\bar x_1,\dots,\bar x_n)=0\iff f\in\mathfrak p.
\end{equation}

\begin{lemma}\label{lem:tangent}
Let $T_y=\bigcap_{f\in\mathfrak p}\ker df_y\subseteq\kappa((t))^n$, let
$\pi\colon\kappa((t))^{n+1}\to\kappa((t))^n$,
$\pi(w)_i=(w_iq_0-q_iw_0)/q_0^2$ $(1\le i\le n)$, whose kernel is the line
$\kappa((t))q$, and put $C=\pi^{-1}(T_y)$.
\begin{enumerate}[label=(\alph*)]
\item $\dim T_y=m:=\dim Y\le d+1$; hence $\dim C=m+1\le d+2$.
\item For every $K$-derivation $\delta$ of $\kappa((t))$, the column
$\delta q=(\delta q_i)_i$ lies in $C$; in particular $q\in C$.
\item If a row vector $\ell\in\kappa((t))^{n+1}$ annihilates $C$, then
$(\ell_1,\dots,\ell_n)\in\operatorname{span}\{df_y:f\in\mathfrak p\}$.
\end{enumerate}
\end{lemma}
\begin{proof}
(a) By the conormal sequence, $\Omega_{F/K}=F^n/\operatorname{span}\{df(\bar x):f\in
\mathfrak p\}$, and $\dim_F\Omega_{F/K}=\operatorname{trdeg}_KF=m$ in characteristic
zero. The rank $n-m$ of the Jacobian is preserved under
$F\subset\kappa((t))$, so $\dim T_y=m$. For the bound, the exact sequence
$\mathfrak m_V/\mathfrak m_V^2\to\kappa\otimes_V\Omega_{V/K}\to\Omega_{\kappa/K}\to0$
gives $\dim_\kappa(\kappa\otimes_V\Omega_{V/K})\le1+d$. By Nakayama (and
Lemma~\ref{lem:place}(e)), $\Omega_{V/K}$ is generated by at most $d+1$ elements,
hence $m=\dim_F(F\otimes_V\Omega_{V/K})\le d+1$. Finally
$\dim C=\dim T_y+1$ since $\pi$ is surjective with one-dimensional kernel.

(b) For $f\in\mathfrak p$ we have $f(y)=0$ by \eqref{eq:generic}, and $\delta$
kills $K$, so $0=\delta(f(y))=\sum_i\partial_if(y)\,\delta y_i$, i.e.\
$\delta y\in T_y$. By the quotient rule, $\pi(\delta q)_i=\delta(q_i/q_0)=\delta y_i$.

(c) We have $\ell(q)=0$ and $\ell(w)=0$ for $w\in C$. Hence for $w\in C$,
\[
 \sum_{i\ge1}\ell_i\pi(w)_i=\frac{q_0\sum_{i\ge1}\ell_iw_i-w_0\sum_{i\ge1}\ell_iq_i}{q_0^2}
 =\frac{-q_0\ell_0w_0+w_0\ell_0q_0}{q_0^2}=0.
\]
As $\pi(C)=T_y$, $(\ell_1,\dots,\ell_n)\in T_y^\perp=\operatorname{span}\{df_y:
f\in\mathfrak p\}$.
\end{proof}

\begin{lemma}\label{lem:lift}
Let $M$ be an $(n+1)\times r$ matrix over $\kappa[[t]]$ whose reduction $M(0)$
has rank $r$, and let $u\in\kappa^{n+1}$ be a row vector with $uM(0)=0$. Then
there is a row vector $\ell(t)\in\kappa[[t]]^{n+1}$ with $\ell M=0$ and
$\ell(0)=u$.
\end{lemma}
\begin{proof}
There is a set $I$ of $r$ row indices with $M_I(0)$ invertible, so $M_I$ is
invertible over $\kappa[[t]]$ (its determinant has nonzero constant term). Let $J$
be the complementary set of indices and define
$\ell_J(t)=u_J$, $\ell_I(t)=-u_JM_JM_I^{-1}$. Then
$\ell M=\ell_IM_I+\ell_JM_J=-u_JM_J+u_JM_J=0$, and $uM(0)=0$ gives
$u_I=-u_JM_J(0)M_I(0)^{-1}$, hence $\ell(0)=u$.
\end{proof}

\begin{lemma}\label{lem:special}
Let $P\in K[\xi_1,\dots,\xi_n]$ vanish on $N(Y)$, and let
$\tilde\ell\in\kappa[[t]]^n$ lie in $\operatorname{span}_{\kappa((t))}\{df_y:f\in
\mathfrak p\}$. Then $P(\tilde\ell(0))=0$.
\end{lemma}
\begin{proof}
Write $\tilde\ell=\sum_{k=1}^s\mu_k\nabla g_k(y)$ with $g_k\in\mathfrak p$ and
$\mu_k\in\kappa((t))$, and put
\[\Phi(x,\mu)=P\Bigl(\sum_{k=1}^s\mu_k\nabla g_k(x)\Bigr)=\sum_\alpha\phi_\alpha(x)\mu^\alpha
\in K[x_1,\dots,x_n,\mu_1,\dots,\mu_s],\]
with $\phi_\alpha\in K[x_1,\dots,x_n]$ and $\mu^\alpha=\mu_1^{\alpha_1}\cdots
\mu_s^{\alpha_s}$. For $y'\in Y_{\mathrm{sm}}$ and every $\mu\in K^s$ we have
$\sum_k\mu_kdg_k(y')\in N^*_{y'}Y\subseteq N(Y)$, so $\Phi(y',\mu)=0$. As $K$ is
infinite, a polynomial in $\mu$ vanishing on all of $K^s$ is zero, so
$\phi_\alpha(y')=0$ for all $\alpha$. Since $Y_{\mathrm{sm}}$ is dense in $Y$,
every $\phi_\alpha$ vanishes on $Y$, i.e.\ $\phi_\alpha\in\mathfrak p=I(Y)$. By
\eqref{eq:generic}, $\phi_\alpha(y)=0$ for all $\alpha$. Hence
$\Phi(y,\mu)=\sum_\alpha\phi_\alpha(y)\mu^\alpha=0$ for all $\mu\in\kappa((t))^s$,
in particular $P(\tilde\ell)=\Phi\bigl(y,(\mu_1,\dots,\mu_s)\bigr)=0$. Now
$\tilde\ell\in\kappa[[t]]^n$ and $P$ has coefficients in $K$, so
$P(\tilde\ell)\in\kappa[[t]]$ has constant term $P(\tilde\ell(0))$, which
therefore vanishes.
\end{proof}

\begin{proof}[Proof of Theorem~\ref{thm:asymptotic-conormal-formal}]
We assume $\bar x_1$ is transcendental over $K$ (the other case was treated
above) and use the notation of the Setup.

\medskip\noindent\textbf{Step 1: the tangent frame.}
We apply $D_k$ coefficientwise to $q_i(t)=\sum_rc_{ir}t^r$ and define
$Z_{ik}(t):=\sum_rD_k(c_{ir})t^r$. Since $D_k(\bar q_{i_l})=\delta_{kl}$, the
reduction $Z(0)$ is the identity on the rows $i_1,\dots,i_d$, and since $q_j=1$ is
constant, $Z_{jk}=0$. Rows $0$ and $1$ of $Z$ have order at least $a$ and $b$,
since $D_k$ preserves $t$-adic order.

Let $\tilde Z$ be the $d\times d$ submatrix of $Z$ consisting of the rows
$i_1,\dots,i_d$ and all columns. Its reduction $\tilde Z(0)$ is the identity
matrix, so $\tilde Z$ is invertible over $\kappa[[t]]$. Hence there exists a
unique vector $\lambda=(\lambda_1,\dots,\lambda_d)^{\mathsf T}\in\kappa[[t]]^d$ with
$\tilde Z\lambda=(q'_{i_1},\dots,q'_{i_d})^{\mathsf T}$. We define $w=q'-Z\lambda$,
so that $w_{i_l}=0$ for $l=1,\dots,d$. Its $0$th entry is
$w_0=q_0'-Z_{0\bullet}\lambda$. Here $q_0'=at^{a-1}u_0+t^au_0'$ has order exactly
$a-1$ (the coefficient $au_0(0)$ of $t^{a-1}$ is nonzero because the
characteristic is zero), while $Z_{0\bullet}\lambda$ has order at least $a$,
because every entry of row $0$ of $Z$ has order at least $a$. Hence
$\operatorname{ord}w_0=a-1$. Let $c$ be the minimal order among the entries of $w$;
then $c\le a-1$, and $\tau:=t^{-c}w$ is a vector over $\kappa[[t]]$ with
$\tau(0)\neq0$ and $\tau_{i_l}=0$. In row $1$, both $q_1'$ and $Z_{1\bullet}\lambda$
have order at least $b-1$, so $\operatorname{ord}w_1\ge b-1$, and
$\operatorname{ord}\tau_1\ge b-1-c\ge b-a>0$. Finally $q_j'=0$ and $Z_{j\bullet}=0$
(as $q_j=1$), so $\tau_j=0$.

\medskip\noindent\textbf{Step 2: the limit plane.}
Let $M=[\,q\mid Z_{\bullet1}\mid\dots\mid Z_{\bullet d}\mid\tau\,]$, an
$(n+1)\times(d+2)$ matrix over $\kappa[[t]]$. Its reduction $M(0)$ has full
column rank: in row $j$ only $q(0)$ is nonzero ($q_j=1$); in rows $i_1,\dots,i_d$
the columns $Z_{\bullet k}(0)$ form an identity block while $\tau$ vanishes there;
and $\tau(0)\neq0$. Since some $(d+2)\times(d+2)$ minor of $M(0)$ is nonzero, the
corresponding minor of $M$ is a unit of $\kappa[[t]]$, so the columns of $M$ are
also linearly independent over $\kappa((t))$. Every column of $M(0)$ has a
vanishing $Z_1$-entry: $q_1(0)=0$ (as $b>0$), the entries $Z_{1k}$ have order at
least $b>0$, and $\tau_1(0)=0$. Thus $L_0=\operatorname{im}M(0)\subseteq\{Z_1=0\}$,
and $\dim L_0=d+2$.

By Lemma~\ref{lem:tangent}(b), the columns $q$, $Z_{\bullet k}$ and $q'$ lie in $C$
(each is $\delta q$ for a $K$-derivation $\delta$ of $\kappa((t))$: the $D_k$
extended coefficientwise, and $d/dt$), hence so does $\tau$, a $\kappa((t))$-linear
combination of them. By Lemma~\ref{lem:tangent}(a), $\dim C\le d+2$, so the $d+2$
independent columns of $M$ span $C$.

\medskip\noindent\textbf{Step 3: lifting and specializing.}
Let $u=(0,1,0,\dots,0)$ be the row vector reading off the $Z_1$-coordinate; since
$L_0\subseteq\{Z_1=0\}$, $uM(0)=0$. By Lemma~\ref{lem:lift} there is
$\ell(t)\in\kappa[[t]]^{n+1}$ with $\ell M=0$ and $\ell(0)=u$. Since the columns of
$M$ span $C$, $\ell$ annihilates $C$, so by Lemma~\ref{lem:tangent}(c) its affine
part $\tilde\ell=(\ell_1,\dots,\ell_n)$ lies in $\operatorname{span}\{df_y:f\in
\mathfrak p\}$, and since $\ell(0)=u$, $\tilde\ell(0)=(1,0,\dots,0)$. Let
$P\in K[\xi_1,\dots,\xi_n]$ vanish on $N(Y)$. By Lemma~\ref{lem:special},
$P(1,0,\dots,0)=0$, which is statement~(i).

For statement~(ii): $\operatorname{Dir}(Y)$ is the common zero set of the
homogeneous polynomials vanishing on $\{[\xi]:0\neq\xi\in N(Y)\}$, equivalently on
$N(Y)$. Each such polynomial vanishes at $(1,0,\dots,0)$ by (i), hence at its
class $[dx_1]=[1:0:\dots:0]$ (homogeneity makes the vanishing independent of the
representative). Therefore $[dx_1]\in\operatorname{Dir}(Y)$. \qed
\end{proof}
\end{comment}
Throughout this section $K$ denotes an algebraically closed field of
characteristic zero. Let $Y\subset X=A^n_{K}$ be an irreducible closed subvariety and $\mathfrak p=I(Y)$ its prime ideal, so $Y=V(\mathfrak p)$.
Because $\mathbb A^n_K$has a global coordinate frame, $T^*X=X\times(K^n)^*$ holds and we can therefore identify all cotangent fibres with $(K^n)^*$. We write $[Z_0:\cdots:Z_n]$ for homogeneous coordinates on the projective space $\mathbb P^n$, and embed $X\hookrightarrow\mathbb P^n$ as the affine chart $\{Z_0\ne0\}$ via
$[1:x_1:\cdots:x_n]$, so that $x_i=Z_i/Z_0$ and the hyperplane at infinity
is $\{Z_0=0\}$.  We define $\operatorname{Dir}(Y)$ to be the Zariski closure in  $\mathbb P((K^n)^*)$  of the nonzero conormal vectors at all smooth affine points of $Y$ projected onto $\mathbb P((K^n)^*)$.

\begin{definition}[Asymptotic conormal directions]\label{def:dir}
For $Y$ as above, the set of \emph{conormal directions} is
\[ \operatorname{Dir}(Y)=\overline{\{[\xi]\in\mathbb P((K^n)^*): 0\neq\xi\in N(Y)\}}
 \subseteq\mathbb P((K^n)^*),\]
the Zariski closure of the projectivization of $N(Y)$.
\end{definition}

Let $\mathfrak p\subset K[x_1,\dots,x_n]$ be prime and define $Y=V(\mathfrak p)$; we assume $Y\cap H = \varnothing$.
We furthermore write $R=K[x]/\mathfrak p$, $F=\operatorname{Frac}(R)$ for the field of fractions, and $\bar x_i$ for the
image of $x_i$ in $R$. 

\begin{lemma}\label{lem:place}
Let $\bar x_1$ be transcendental over $K$. 
Put $Z_0=1$, $Z_i=\bar x_i$ $(1\le i\le n)$. There exist an index $j$, an element $s\in F^\times$, a
finitely generated integrally closed $K$-subalgebra $A\subseteq F$ with $\operatorname{Frac}A=F$, and a height-one
prime $\mathfrak p_D\subset A$ such that, with $q_i:=Z_is$ and $V:=A_{\mathfrak p_D}$:
\begin{enumerate}[label=(\alph*)]
\item $q_j=1$, $q_i\in A$ for all $i$, $\bar x_1\in\mathfrak p_D$, and $A$ is integral over $C_0:=K[q_0,\dots,q_n,\bar x_1]$;
\item $V$ is a DVR with $K\subseteq A\subseteq V$, $\operatorname{Frac}V=F$, $\bar x_1\in\mathfrak m_V$, and $R\not\subseteq V$.
\end{enumerate}
\end{lemma}
\begin{proof}
\emph{(1) A valuation ring.} Since $\bar x_1$ is transcendental over $K$, $S=K[\bar x_1]$ is a polynomial ring and
$\bar x_1S$ is a proper ideal. By Chevalley's extension theorem there is a valuation ring $W\subseteq F$ with
$S\subseteq W$ and $\bar x_1\in\mathfrak m_W$. In particular $K\subseteq W$ and $\bar x_1^{-1}\notin W$.

\emph{(2) A chart.} Among the nonzero $Z_a$ choose $j$ such that $Z_a/Z_j\in W$ for all $a$ (possible since
$W$ is a valuation ring and the index set is finite), and put $s=Z_j^{-1}$. Then $q_j=1$ and $q_a\in W$ for all $a$.
Since $\bar x_i=q_i/q_0$ and $q_0=s\neq0$, the ring $C_0=K[q_0,\dots,q_n,\bar x_1]\subseteq W$ is finitely generated with
$\operatorname{Frac}C_0=F$.

\emph{(3) Normalization.} Let $A$ be the integral closure of $C_0$ in $F$. It is finite over $C_0$ (finiteness of
normalization for finitely generated domains over a field), hence a finitely generated, integrally closed,
Noetherian domain with $\operatorname{Frac}A=F$. Since $W$ is integrally closed, $A\subseteq W$, so
$\bar x_1^{-1}\notin A$ and $\bar x_1$ is a nonzero nonunit of $A$.

\emph{(4) The divisor.} By Krull's principal ideal theorem there is a prime $\mathfrak p_D\ni\bar x_1$ of $A$ of height
one. As $A$ is normal and Noetherian, $V=A_{\mathfrak p_D}$ is a DVR with $\operatorname{Frac}V=F$, $A\subseteq V$ and
$\bar x_1\in\mathfrak m_V$. (Geometrically, $V$ is the local ring of a prime divisor of the normalization of the
affine chart $\{Z_j\ne0\}$ of $\overline Y$. No properness is used.)

\emph{(5)} If $R\subseteq V$, then $\bar x_1$ would be a unit of $V$ because $\bar x_1^{-1}\in R$, contradicting $\bar x_1\in\mathfrak m_V$.
\end{proof}



\begin{lemma}[Residue field]\label{lem:kappa}
Let $\kappa=V/\mathfrak m_V=\operatorname{Frac}(A/\mathfrak p_D)$ and $\bar q_i$ the residue of $q_i$. Then:
\begin{enumerate}[label=(\alph*)]
\item $\kappa$ is algebraic over $L=K(\bar q_0,\dots,\bar q_n)$.
\item $d:=\dim_\kappa\Omega_{\kappa/K}$ is finite, and there are indices $i_1,\dots,i_d\notin\{0,1,j\}$ such that
$d\bar q_{i_1},\dots,d\bar q_{i_d}$ form a basis of $\Omega_{\kappa/K}$. Let $D_1,\dots,D_d$ be the dual
$K$-derivations of $\kappa$, so $D_k(\bar q_{i_l})=\delta_{kl}$.
\item $\Omega_{V/K}$ is a finitely generated $V$-module.
\end{enumerate}
\end{lemma}
\begin{proof}
(a) Let $a'\in A$. As $A$ is integral over $C_0=K[q,\bar x_1]$, an integral equation of $a'$ over $C_0$
reduces modulo $\mathfrak m_V$ to an integral equation of $\bar a'$ over the image of $C_0$ in $\kappa$. This image
is $K[\bar q]$, because $\bar x_1\in\mathfrak m_V$, so $\bar a'$ is algebraic over $L$. Every element of $\kappa$ has the form
$\bar a'/\bar b'$ with $a',b'\in A$, $b'\notin\mathfrak p_D$, and algebraic elements form a field.

(b) $\kappa$ is a finitely generated field extension of $K$. Since $\bar q_0=\bar q_1=0$ and $\bar q_j=1$, the field $L$ is generated by the
remaining $\bar q_i$. Choose a transcendence basis $B=\{\bar q_{i_1},\dots,\bar q_{i_d}\}$ of $L/K$ among them; then
$i_l\notin\{0,1,j\}$. By (a), $\kappa/K(B)$ is algebraic, hence separable in characteristic $0$, so
$\Omega_{\kappa/K(B)}=0$ and $d\bar q_{i_1},\dots,d\bar q_{i_d}$ is a basis of $\Omega_{\kappa/K}$. The dual basis of
$\operatorname{Hom}_\kappa(\Omega_{\kappa/K},\kappa)=\operatorname{Der}_K(\kappa,\kappa)$ gives the $D_k$.

(c) $A$ is a finitely generated $K$-algebra, so $\Omega_{A/K}$ is a finite $A$-module, and
$\Omega_{V/K}=V\otimes_A\Omega_{A/K}$ because $V=A_{\mathfrak p_D}$ is a localization.
\end{proof}

\begin{lemma}[Tangent space at $y$]\label{lem:tangent}
Let $T_y=\bigcap_{f\in\mathfrak p}\ker df_y\subseteq\kappa((t))^n$ and let
$\pi\colon\kappa((t))^{n+1}\to\kappa((t))^n$,
$\pi(w)_i=(w_iq_0-q_iw_0)/q_0^2$ ($1\le i\le n$), with kernel $\kappa((t))q$. Put
$C=\pi^{-1}(T_y)$.
\begin{enumerate}[label=(\alph*)]
\item $dim T_y = m\leq d+1$; hence $\dim C=m+1\le d+2$.
\item For every $K$-derivation $\delta$ of $\kappa((t))$ the column $\delta q=(\delta q_i)_i$ lies in $C$ (in particular $q\in C$).
\item If a row vector $\ell\in\kappa((t))^{n+1}$ annihilates $C$, then $(\ell_1,\dots,\ell_n)\in\operatorname{span}\{df_y:f\in\mathfrak p\}$.
\end{enumerate}
\end{lemma}
\begin{proof}
(a) By the conormal sequence $\Omega_{F/K}=F^n/\operatorname{span}\{df(\bar x)\}$, and
$\dim_F\Omega_{F/K}=\operatorname{trdeg}_KF=m$ in characteristic $0$
\cite[II.8]{hartshorne}. Rank of the Jacobian is preserved under $F\subset\kappa((t))$, so
$\dim\operatorname{span}\{df_y\}=n-m$ and $\dim T_y=m$. For the bound: $\Omega_{V/K}$ is
a finitely generated $V$-module ($V=A_{\mathfrak p_D}$), and the sequence
$\mathfrak m_V/\mathfrak m_V^2\to\kappa\otimes\Omega_{V/K}\to\Omega_{\kappa/K}\to0$ gives
$\dim_\kappa(\kappa\otimes\Omega_{V/K})\le1+d$. By Nakayama $\Omega_{V/K}$ is generated by
$\le d+1$ elements, hence $m=\dim_F(F\otimes\Omega_{V/K})\le d+1$.
Since $\ker\pi$ is a line, $\dim C=\dim T_y+1$.

(b) For $f\in\mathfrak p$, $f(y)=0$ and $\delta$ kills $K$, so $0=\delta f(y)=\sum_i\partial_if(y)\,\delta y_i$,
i.e.\ $\delta y\in T_y$. The quotient rule gives $\pi(\delta q)_i=\delta(q_i/q_0)=\delta y_i$.

(c) $\ell(q)=0$ and $\ell(w)=0$ for $w\in C$. For $w\in C$,
$\sum_{i\ge1}\ell_i\pi(w)_i=\bigl[q_0\sum_{i\ge1}\ell_iw_i-w_0\sum_{i\ge1}\ell_iq_i\bigr]/q_0^2
=\bigl[-q_0\ell_0w_0+w_0\ell_0q_0\bigr]/q_0^2=0$. As $\pi(C)=T_y$, $(\ell_1,\dots,\ell_n)\in T_y^\perp=\operatorname{span}\{df_y\}$.
\end{proof}





\begin{theorem}\label{thm:asymptotic-conormal-formal}
Let $Y\subseteq\mathbb A^n$ be nonempty, closed, and irreducible with $Y\cap H=\varnothing$, and let
$N(Y)=\bigcup_{y\in Y_{\mathrm{sm}}}N^*_yY\subseteq K^n$ be the conormal cone over the smooth locus $Y_{sm}$.
\begin{enumerate}[label=(\roman*)]
\item Every polynomial $P\in K[\xi_1,\dots,\xi_n]$ vanishing on $N(Y)$ satisfies $P(1,0,\dots,0)=0$.
\item $[dx_1]\in\operatorname{Dir}(Y)$.
\end{enumerate}
\end{theorem}

\AIadd{\formalref[{over every algebraically closed field of characteristic zero; the formal proof follows a different route, written out after the printed proof}]{%
\leandecl{Stafford38/Geometry/GeneralAsymptoticConormal.lean}{42}{Stafford38.Geometry.GeneralAsymptoticConormal.coordinate_axis_mem_projective_conormal_directions},
affine-cone form \leandecl{Stafford38/Geometry/GeneralAsymptoticConormal.lean}{32}{Stafford38.Geometry.GeneralAsymptoticConormal.coordinate_axis_mem_smooth_fibre_closure}}}
\begin{proof}
If $\bar x_1$ is algebraic over $K$, there is a nonzero polynomial $g$ over $K$
in one variable with $g(\bar x_1)=0$, i.e.\ $g(x_1)\in\mathfrak p$. Since $K$
is algebraically closed, $g(x_1)=\lambda\prod_i(x_1-c_i)$ with
$\lambda\in K^\times$, and as $\mathfrak p$ is prime, some factor $x_1-c$ lies
in $\mathfrak p$. Thus $x_1=c$ on $Y$, and $c\neq0$ because $Y\neq\varnothing$
avoids $x_1=0$. In this case $f=x_1-c\in\mathfrak p=I(Y)$ has differential
$df_y=dx_1$ at every point $y\in Y_{\mathrm{sm}}$, so $dx_1\in N_y^*Y$ and its coordinate vector in the dual basis $(1,0,\dots , 0) \in N(Y)$. It follows directly that every polynomial vanishing on $N(Y)$ therefore vanishes at $(1,0,\dots , 0)$, and
$[dx_1]\in\operatorname{Dir}(Y)$. So we may assume from now on that
$\bar x_1$ is transcendental over $K$.\\

The rest of the proof follows this rough logic: We build an arc to a place at
infinity where $x_1$ vanishes, with generic point $y$ on $Y$. We show that the
tangent planes $L(t)$ to $Y$ at $y$, $t\neq0$, have a limit plane $L_0$ of
the same dimension at $t=0$, and that $L_0\subseteq\{Z_1=0\}$. Hence the
covector $Z_1^\vee$ reading off the $Z_1$-coordinate annihilates $L_0$, and we
can lift it to a family of covectors $\ell(t)$ annihilating $L(t)$ for
$t\neq0$, with $\ell(0)=Z_1^\vee$. For $t\neq0$ the $\ell(t)$ are conormal
covectors to $Y$, so their limit $\ell(0)$, whose affine part is $dx_1$, lies
in the closure of the conormal directions, i.e.\ $[dx_1]\in\operatorname{Dir}(Y)$.\\


Since $\bar x_1$ is transcendental over $K$, by Lemma~\ref{lem:place} there exists a divisorial valuation ring $V\subset F$
with $\bar x_1\in\mathfrak m_V$ (the maximal ideal) and $R\not\subseteq V$. Let $\varpi$ be a
generator of $\mathfrak m_V$, then every nonzero $f \in F$ can be written as $f = \varpi^k\cdot u$ with $k \in \mathbb{Z}$ and $u$ a unit of $V$ - we say $\operatorname{ord}_V(f) = k \in \mathbb Z$ is the valuation of $f$ in $V$.  We denote the corresponding residue field by $\kappa=V/\mathfrak m_V$. 

Let $\widehat V$ be the completion of $V$. Since $V\supseteq K$
and $K$ is of characteristic zero, there exists a coefficient field of $\widehat V$
containing $K$ and mapping isomorphically onto $\kappa$ \cite[Tag~0C0S]{Stacks};
choosing one gives use an isomorphism $\widehat V\cong\kappa[[t]]$ under which
$\varpi$ corresponds to the formal variable $t$. Since the natural map $V\to\widehat V$ is injective ($\bigcap_k\varpi^kV=0$) and extends to an
embedding $F\hookrightarrow\kappa((t))$, we identify elements of $F$ with
Laurent series.

Let $\bar x_0:=1$. Choose $j\in\{0,\dots,n\}$ and a valuation ring $W$ such that
$\bar x_a/\bar x_j\in W$ for all $a$, and define $q_i:=\bar x_i/\bar x_j$ for
$0\le i\le n$.
Set $Z_0=1$, $Z_i=\bar x_i$ and let
$j$ and $q=(q_0,\dots,q_n)$ be as in Lemma~\ref{lem:place}, so that
$q_i\in V$, $q_j=1$, $q_1=q_0\bar x_1$ (since $q_1/q_0 = Z_1/Z_0 = \bar x_1$). Let $a=\operatorname{ord}_V(q_0)$ and $b=\operatorname{ord}_V(q_1)$.  $a,b\geq 0$ since $q_0,q_1$ are elements of $V$, so that in $\kappa[[t]]$
\[q_0=t^au_0,\qquad q_1=t^bu_1,\]
with $u_0(0)\neq0\neq u_1(0)$. Since $\bar x_1\in\mathfrak m_V$, $b-a=\operatorname{ord}_V(q_1/q_0) = \operatorname{ord}_V(\bar x_1)>0$.  Moreover
$a>0$: if $a=0$, then $q_0$ is a unit of $V$, hence all $\bar x_i=q_i/q_0$
lie in $V$, so $R\subseteq V$ since also $K\subseteq V$, contradicting Lemma~\ref{lem:place}(b). It follows that $j \notin \{0,1\}$ since $q_j = 1$ is a unit while $q_0,q_1$ are not.

Dividing by $q_0\neq 0$, we define $y=(q_1/q_0,\dots,q_n/q_0)\in\kappa((t))^n$; its
coordinates are the images of $\bar x_1,\dots,\bar x_n$ under $F\hookrightarrow
\kappa((t))$. Since ring homomorphisms commute with polynomial evaluation and
this map is injective, for every polynomial $f\in K[x_1,\dots,x_n]$ we have
\begin{equation}
    f(y)=0\iff f(\bar x_1,\dots,\bar x_n)=0 \iff  f\in\mathfrak p.
    \label{eq:generic}
\end{equation}
$f(y)=0$ if and only if $f(\bar x_1,\dots,\bar x_n)=0$ in $F$, i.e.\ if and
only if $f\in\mathfrak p$. \\


We will now build independent vectors from the tangential planes of $q$.
Let $d= \dim_\kappa \Omega_{\kappa/K}$, and let $i_1,\dots,i_d\notin\{0,1,j\}$ and
$D_1,\dots,D_d$ be as in Lemma~\ref{lem:kappa}. 

We apply $D_k$ coefficientwise to $q_i(t)=\sum_rc_{ir}t^r$ and define $Z_{ik}(t):=\sum_rD_k(c_{ir})t^r$. Since $D_k(\bar q_{i_l})=\delta_{kl}$, the reduction $Z(0)$ is the identity on the rows $i_1,\dots,i_d$, and since $q_j = 1$ is constant, $Z_{jk}=0$. Rows $0$ and $1$ of $Z$ have order at least $a$ and $b$, since $D_k$ preserves $t$-adic order.

Let $\tilde Z$ be the $d\times d$ submatrix of $Z$ consisting of the rows
$i_1,\dots,i_d$ and all columns. Its reduction $\tilde Z(0)$ is the identity
matrix, so $\tilde Z$ is invertible over $\kappa[[t]]$. Hence there exists a unique vector
$\lambda=(\lambda_1,\dots,\lambda_d)^{\mathsf T}\in\kappa[[t]]^d$ with
$\tilde Z\lambda=(q'_{i_1},\dots,q'_{i_d})^{\mathsf T}$. We define $w=q'-Z\lambda$, so
that $w_{i_l}=0$ for $l=1,\dots,d$.
Its $0$th entry is $w_0=q_0'-Z_{0\bullet}\lambda$. Here
$q_0'=at^{a-1}u_0+t^au_0'$ has order exactly $a-1$ (since the coefficient
$au_0(0)$ of $t^{a-1}$ is nonzero),
while $Z_{0\bullet}\lambda$ has order at least $a$, because every entry of
row $0$ of $Z$ has order at least $a$ (since $q_0$
has order $a$). Hence $\operatorname{ord}w_0=a-1$. Let $c$ be the minimal order
among the entries of $w$; then $c\le a-1$, and $\tau:=t^{-c}w$ is a vector
over $\kappa[[t]]$ with $\tau(0)\neq0$ and $\tau_{i_l}=0$. In row $1$, both
$q_1'$ and $Z_{1\bullet}\lambda$ have order at least $b-1$, so $\operatorname{ord}w_1\ge b-1$, so
$\operatorname{ord}\tau_1\ge b-1-c\ge b-a>0$. Finally $q_j'=0$ and
$Z_{j\bullet}=0$ (as $q_j=1$), so $\tau_j=0$. \\


We can show that $q, Z_{\bullet1},\dots, Z_{\bullet d}, \tau$ span all of $C$: 
Let $M=[\,q\mid Z_{\bullet1}\mid\dots\mid Z_{\bullet d}\mid\tau\,]$, an
$(n+1)\times(d+2)$ matrix over $\kappa[[t]]$. Its reduction $M(0)$ has full
column rank: in row $j$ only $q(0)$ is nonzero ($q_j=1$); in rows $i_1,\dots,i_d$
the columns $Z_{\bullet k}(0)$ form an identity block while $\tau$ vanishes
there; and $\tau(0)\neq0$, i.e. we have $d+2$ independent reduced column vectors. This also translates to the column vectors of $M$ themselves: Since some $(d+2)\times(d+2)$ minor of $M(0)$ is nonzero, the corresponding
minor of $M$ is a unit of $\kappa[[t]]$; hence the columns of $M$ are also
linearly independent over $\kappa((t))$, so they span a $(d+2)$-dimensional
subspace there.
Furthermore, every column of $M(0)$ has a vanishing $Z_1$-entry:
$q_1(0)=0$ (as $b>0$), the entries $Z_{1k}$ have order at least $b>0$, and
$\tau_1(0)=0$. Thus the span $L_0$ of $q(0), Z_{\bullet1}(0),\dots, Z_{\bullet d}(0), \tau(0)$ is contained in $\{Z_1=0\}$ and also has dimension $d+2$. 

By Lemma~\ref{lem:tangent}(b), the columns $q$, $Z_{\bullet k}$ and
$q'$ lie in $C$ (each is $\delta q$ for a $K$-derivation $\delta$ of
$\kappa((t))$: the $D_k$ extended coefficientwise, and $d/dt$), hence so does $\tau$,
a combination of them. By Lemma~\ref{lem:tangent}(a), $\dim C\le d+2$, so the
$d+2$ independent columns of $M$ must span $C$.\\


Let $u=(0,1,0,\dots,0)$ be the row vector reading off the $Z_1$-coordinate;
since $L_0=\operatorname{im}M(0)\subseteq\{Z_1=0\}$ by Step~3, $uM(0)=0$.
We can lift $u$ to a row vector $\ell(t)\in\kappa[[t]]^{n+1}$ with $\ell M=0$ and
$\ell(0)=u$: Since $M(0)$ has full column rank $d+2$, there is a set $I$ of
$d+2$ row indices for which the square submatrix $M_I(0)$ of $M(0)$ is
invertible; then $M_I$ is invertible over $\kappa[[t]]$, since its determinant
has a nonzero constant term. Let $J$ be the complementary set of row indices,
and define $\ell$ by
\[\ell_J(t)=u_J,\qquad \ell_I(t)=-u_JM_JM_I^{-1},\]
where $M_J$ denotes the submatrix of $M$ with rows in $J$. Then
$\ell(t) M=\ell_I(t)M_I+\ell_J(t)M_J=-u_JM_J+u_JM_J=0$. Moreover
$uM(0)=0$ gives $u_I=-u_JM_J(0)M_I(0)^{-1}$, hence $\ell(0)=u$.

Since the columns of $M$ span $C$, $\ell$ annihilates $C$, so by
Lemma~\ref{lem:tangent}(c) its affine part
$\tilde\ell(t)=(\ell_1(t),\dots,\ell_n(t))$ lies in
$\operatorname{span}\{df_y:f\in\mathfrak p\}$, and since $\ell(0)=u$,
$\tilde\ell(0)=(1,0,\dots,0)$.

Finally, let $P\in K[\xi_1,\dots,\xi_n]$ vanish on $N(Y)$. We will show
$P(\tilde\ell(0))=0$: Since $\tilde\ell\in\operatorname{span}_{\kappa((t))}\{dg_y:g\in
\mathfrak p\}$, we can write
\[\tilde\ell=\sum_{k=1}^s\mu_k\nabla g_k(y),\qquad g_k\in\mathfrak p,\
 \mu_k\in\kappa((t)),\]
 With some $s \in \mathbb{N}$. Here $\nabla g(x)=(\partial g/\partial x_1(x),\dots,
\partial g/\partial x_n(x))$ denotes the coordinate vector of the covector
$dg_x$. We furthermore define
\[\Phi(x,\mu)=P\Bigl(\sum_{k=1}^s\mu_k\nabla g_k(x)\Bigr)
 =\sum_\alpha\phi_\alpha(x)\mu^\alpha\in K[x_1,\dots,x_n,\mu_1,\dots,\mu_s],\]
with $\phi_\alpha\in K[x_1,\dots,x_n]$ and $\mu^\alpha = \mu_1^{\alpha_1}\dots\mu_s^{\alpha_s}$ for a multi-infex $\alpha \in \mathbb{N}^s$. 
For a smooth point $y'\in Y_{\mathrm{sm}}$ and
$\mu\in K^s$ we have $\sum_k\mu_kdg_k(y')\in N^*_{y'}Y\subseteq N(Y)$, so
$\Phi(y',\mu)=0$. As $K$ is infinite, a polynomial in $\mu$ vanishing on all of
$K^s$ is zero, so $\phi_\alpha(y')=0$ for all $\alpha$. Since $Y_{\mathrm{sm}}$
is dense in $Y$, every $\phi_\alpha$ vanishes on all of $Y$, i.e.\
$\phi_\alpha\in\mathfrak p=I(Y)$. By \eqref{eq:generic} (applied to the generic
point $y\in\kappa((t))^n$), we therefore have $\phi_\alpha(y)=0$ for all
$\alpha$. Hence $\Phi(y,\mu)=\sum_{\alpha}\phi_\alpha(y)\mu^\alpha=\sum_\alpha 0\cdot \mu^\alpha = 0$ for all $\mu\in\kappa((t))^s$, and in
particular
\[ P(\tilde\ell)=\Phi\bigl(y,(\mu_1,\dots,\mu_s)\bigr)=0.\]
Now $\tilde\ell\in\kappa[[t]]^n$ and $P$ has coefficients in $K$, so
$P(\tilde\ell)\in\kappa[[t]]$, and its constant term is $P(\tilde\ell(0))$. Since $P(\tilde\ell)=0$, we get
$P(\tilde\ell(0))=P(1,0,\dots,0)=0$, which is statement~(i).

We have shown that every polynomial $P$ vanishing on $N(Y)$ vanishes at
$(1,0,\dots,0)$ (statement~(i)). By the definition of the Zariski closure,
$\operatorname{Dir}(Y)$ is the common zero set of all homogeneous polynomials
vanishing on $\{[\xi]:0\ne\xi\in N(Y)\}$, equivalently on $N(Y)$. Each such
polynomial vanishes at $(1,0,\dots,0)$, hence at its class
$[dx_1]=[1:0:\dots:0]$ (homogeneity makes the vanishing independent of the
representative). Therefore $[dx_1]\in\operatorname{Dir}(Y)$. 
\end{proof}


\begin{comment}
\begin{proof}
If $Y$ is zero-dimensional, its tangent space is trivial, so every cotangent direction occurs in its conormal
fibre. If $\dim Y>0$, but the coordinate function $x_1|_Y$ is constant on $Y$, it is a nonzero constant (since $Y\cap \{x_1=0\} = \varnothing$) and $dx_1$ annihilates every tangent space and is therefore a conormal vector to $Y$, i.e. $[dx_1]\in \operatorname{Dir}(Y)$. We may therefore assume that $m=\dim Y>0$ and that $x_1|_Y$ is nonconstant.

Let $\overline Y\subset\mathbb P^n$ be the closure of $Y$ in $\mathbb P^n$ and let
$\nu:\widetilde Y\to\overline Y$ be its normalization. Normalization is
finite in this setting \cite[\AIreplace{Tag 0BXQ}{Tag~0BXR}, Lemma~33.27.1]{Stacks}.
The rational function $x_1=Z_1/Z_0$ has a zero prime divisor $D$\AIadd{ }on
$\widetilde Y$\AIreplace{, ,}{,} i.e. a prime divisor with $\ord_D(x_1)>0$. Otherwise $x_1^{-1}$ would have no pole at any codimension-one point of $\widetilde Y$. Since $\widetilde Y$ is normal,
$x_1^{-1}$ would be globally regular \cite[Tag 031T]{Stacks}.
As $\widetilde Y$ is finite (hence proper) over $\overline Y$, it is proper over $\mathbb C$.
\AIadd{Finite normalization and preservation of projectivity are standard\cite[\AIreplace{Tag~0BXQ}{Tags~0BXR, 0GK4}, Lemmas~33.27.1--2]{Stacks}.} By \cite[\AIremove{Tag 04L0}\AIadd{Tag 04L2}, Lemma~33.26.2]{Stacks}, any globally regular function must be constant. This contradicts our standing assumption that $x_1|_Y$ is nonconstant.
Since $Y\cap \{x_1=0\}=\varnothing$, $x_1$ \AIreplace{does not vanish on $Y$}{is a unit on $Y$ (Nullstellensatz)}, so the generic point
of $D$ lies outside $\nu^{-1}(Y)$, hence in the hyperplane $\{Z_0 = 0\}$ at infinity. %Because we defined $Y$ to be closed in $X$, $\overline{Y}\backslash Y$ cannot lie in $X$, hence at infinity, $\nu^{-1}(\overline{Y}\backslash Y)$. %above the hyperplane
%The rational function \[ x_1=Z_1/Z_0\] has a zero prime divisor $D$ on $\widetilde Y$, i.e. an irreducible closed subvariety of codimension 1 such that $x_1$ vanishes on $D$. Otherwise $x_1^{-1}$ would have no pole at any codimension-one point. Normality would make it globally regular \cite[Tag 031T]{Stacks}, and projectivity would make it constant \cite[Tag 04L0, Lemma~33.26.2]{Stacks}.
We choose $Z_j$ nonzero at the generic point of $D$, and define
$q_i=Z_i/Z_j$. With respect to the divisorial valuation at $D$, we define
\begin{equation}a=\operatorname{ord}_D(q_0),\qquad
 b=\operatorname{ord}_D(q_1).
\label{eq:ab}\end{equation}
Indeed, because the generic point lies in $\{Z_0=0\}$, $a$ must be positive\AIadd{; as $a,b>0$ while $q_j=1$, the index $j$ differs from $0$ and $1$}. $b$ is finite since $Z_1$
does not vanish identically on $\overline Y$, and we recall $x_1=Z_1/Z_0=q_1/q_0$, therefore $\ord_D(x_1) = \ord_D(q_1)-\ord_D(q_0) = b-a$. Because $D$ was specifically chosen to be a zero prime divisor of $x_1$, $\ord_D(x_1)>0$, hence $b>a.$
We shrink a neighbourhood $U$ of the generic point of $D$ so that $Z_j$ is
invertible. Defining $D\cap U=(s)$, we can find regular functions $w_0$ and $w_1$ that do not vanish identically on $D$, such that
\[ q_0=s^aw_0,\qquad q_1=s^bw_1,\]
Restricting the normalization to $D$ inherits finiteness, so at the generic point of $D$, the map has finite fibers.
Because $\nu_{|D}$ is finite, its function-field extension is separable since $\mathbb C$ is a field of characteristic 0 (\cite[Tag 0322]{Stacks}), and consequently, $\nu_{|D}$ is smooth at the generic point of $D$ \cite[Tag 07ND]{Stacks}\AIadd{, applied to the dominant finite morphism $D\to\nu(D)$ onto its image}.
 Since smoothness is an open condition, $\nu_{|D}$
is smooth on a dense open subset of $D$. Equivalently, since $D$ and $\nu(D)$ have equal dimension
$m-1$, its differential has rank $m-1$.
\AIcomment{GEO-01}{\textbf{Proposed change --- author acceptance pending.} The generic-smoothness step is valid for the dominant
finite map $D\to\nu(D)$ (Tag~07ND). Name that image as the target and
choose the point also over its smooth locus; then the differential into
the projective chart has the asserted rank $m-1$.}

Choose a closed point $p$ of $D\cap U$ in the smooth loci of both $D$ and
$\widetilde Y$, where the differential of $\nu_{|D}$ has rank $m-1$, and
outside $V(w_0w_1)$, i.e. $w_0\neq 0$ and $w_1\neq 0$ at the chosen point\AIadd{, and lying over the smooth locus of $\nu(D)$}. Normality of $\widetilde Y$ ensures that its singular locus does not contain $D$ and because all conditions hold on a dense open subset of $D$, their finite intersection is nonempty, so such a point exists.

%tag 0C0S Lemma 10.160.10 says that completed local ring is isomorphic to a power series ring over its residue field
We can define a completed local ring at $p$, which is isomorphic to a power series ring \cite[Tag 0C0S, Lemma 10.160.10]{Stacks}. We can choose a formal coordinate system centered around $p$ with parameters
$(t,z_2,\ldots,z_m)$. We choose $t=s$: since our point lies in the smooth locus of $D = \{s = 0\}$, $ds\neq 0$ there, so $s$ extends to a regular system of parameters.
\AIadd{More specifically, choose distinct rows $i_1,\ldots,i_{m-1}$,
away from $0,1,j$, such that the restrictions of $dq_{i_k}$ to $T_pD$
are independent; these exist by the rank condition. Take
$z_\alpha=q_{i_{\alpha-1}}-q_{i_{\alpha-1}}(p)$ for $\alpha=2,\ldots,m$.
Together with $t=s$ these form local parameters, because their restrictions
form a cotangent basis on $D$ and $ds$ is transverse to $D$.}


The set $\widetilde Y\setminus\nu^{-1}(Y_{\mathrm{sm}})$ (with $Y_{\mathrm{sm}}$ being the smooth locus of $Y$) is a proper closed subset of $\widetilde Y$ containing $p$. Let $I\subset
\widehat{\mathcal O}_{\widetilde Y,p}\cong\mathbb C[[t,z_2,\ldots,z_m]]$ be
the ideal of this closed set in the completed local ring, so that every element of $I$ vanishes at every point of this closed set, and choose any
nonzero $f\in I$. We expand $f$ in homogeneous parts,
\[ f = f_d(t,z_2,\ldots,z_m) + f_{d+1}(t,z_2,\ldots,z_m) + \cdots,\]
where $f_d\ne0$ is the lowest-degree nonzero homogeneous part. Since
$f_d(1,z_2,\ldots,z_m)$ is a nonzero polynomial in $z_2,\ldots,z_m$, it
vanishes only on a proper closed subset of $\mathbb C^{m-1}$. We now choose
constants $c_2,\ldots,c_m$ outside this subset, so that
\[f_d(1,c_2,\ldots,c_m)\ne0.\] When we change the coordinate system by replacing $z_\alpha$ by
$z_\alpha-c_\alpha t$, we therefore guarantee that $f\ne0$ for $t\ne0$
along the (tilted) formal $t$-axis. Since $f$ vanishes on all of
$\widetilde Y\setminus\nu^{-1}(Y_{\mathrm{sm}})$ (as $f$ lies in its
ideal), this forces every point of the arc with $t\ne0$ -- in particular
its generic point -- to lie outside $\widetilde Y\setminus
\nu^{-1}(Y_{\mathrm{sm}})$, i.e.\ in $\nu^{-1}(Y_{\mathrm{sm}})$.

Setting $z_2=\cdots=z_m=0$ in these coordinates therefore defines a
transverse map $\gamma:\operatorname{Spec}\mathbb C[[t]]\to\widetilde Y$
whose generic point lies, by construction, in $Y_{\mathrm{sm}}$, the smooth
affine locus of $Y$.

In the new coordinates of the local ring, we can accordingly rewrite
\begin{equation} q_0=t^au_0,\qquad q_1=t^bu_1,
 \label{eq:q-orders}\end{equation}
where $u_0,u_1$ are units since $w_0(p)\neq 0, w_1(p)\neq 0$.

In the coordinates $(t,z_2,\ldots,z_m)$, each $q_i$ is a regular function
on $\widetilde Y$ near $p$, so $v=(q_0,\ldots,q_n)^\top$, with $q_j = 1$  is a local
parametrization, in these coordinates, of $\nu$'s image in the chart
$Z_j\ne0$.
We define the partial derivative of $v$ with respect to $t$ and the $z_\alpha$
\[ V_\alpha=\frac{\partial v}{\partial z_\alpha},\qquad
 T=\frac{\partial v}{\partial t}. \]
$V_\alpha$ vanish in rows $0$, $1$, and $j$ at $t = 0$, since $q_j = 1$ is constant, and $q_0$ and $q_1$ (and hence their derivatives) have factors $t$ vanishing at $t=0$. Since locally, $D = \{t=0\}$, they are
the images of a tangent basis of $D$ under the differential of $\nu_{|D}.$ Because we chose $p$ such that $\nu_{|D}$ would have full rank and is hence injective, the images $V_\alpha$ of the tangent basis of $D$ are also independent.
Since the reductions of $V_2,\ldots,V_m$ are independent and vanish in
rows $0,1,j$, they have rank $m-1$, so there exist $m-1$ further
coordinate rows $i_1,\ldots,i_{m-1}$ (necessarily distinct from $0,1,j$) in
which the corresponding $(m-1)\times(m-1)$ submatrix of $(V_2,\ldots,V_m)$
is invertible. \AIreplace{Because this minor is invertible at $t=0$, it remains
invertible over $\mathbb C[[t]]$, so there are unique coefficients
$\lambda_2(t),\ldots,\lambda_m(t)\in\mathbb C[[t]]$ solving}{With our choice of tangential coordinates, the tilted rows satisfy
$q_{i_k}=q_{i_k}(p)+z_{k+1}+c_{k+1}t$. Thus the transverse minor is the
identity and the correction coefficients are the constants $c_\alpha$:}
\[\AIremove{\sum_{\alpha=2}^m \lambda_\alpha(t)\,(V_\alpha)_{i_k} = T_{i_k},
 \qquad k=1,\ldots,m-1.}\]
\AIadd{\[(V_\alpha)_{i_k}=\delta_{\alpha,k+1},\qquad
 T_{i_k}=c_{k+1},\qquad k=1,\ldots,m-1.\]}
Setting
\[ \AIremove{T' := T - \sum_{\alpha=2}^m \lambda_\alpha(t)\,V_\alpha}\]
\AIadd{\[T' := T - \sum_{\alpha=2}^m c_\alpha\,V_\alpha.\]}
therefore gives a column vector whose entries in rows
$i_1,\ldots,i_{m-1}$ vanish identically by construction.
\AIcomment{GEO-06}{\textbf{Proposed local simplification --- author acceptance pending.}
Choose the selected projective rows as tangential coordinates before tilting.
Then the displayed minor is the identity and $\lambda_\alpha(t)=c_\alpha$;
no power-series matrix inversion is needed. The original wording remains
visible above. The order estimates and tangent-lattice argument continue
unchanged; ROUTE-01 records the current formal correspondence status.}
 The entries of $V_\alpha$ in rows $0$ and $1$ have valuation at least $a$ and $b$ respectively, whereas
\[ \partial_tq_0=a t^{a-1}u_0+O(t^a),\qquad
 \partial_tq_1=b t^{b-1}u_1+O(t^b).\]
Because we work over a field of characteristic zero, the displayed
leading coefficients are nonzero, so
\begin{equation} \operatorname{ord}_t(T_0)=a-1,\qquad \operatorname{ord}_t(T_1)=b-1.
 \label{eq:T-orders-raw}\end{equation}
Since the entries of $V_\alpha$ in rows $0,1$ have strictly higher
valuation, subtracting a combination of $V_2,\ldots,V_m$ from $T$ to form
$T'$ cannot change these orders, hence
\begin{equation} \operatorname{ord}_t(T'_0)=a-1,\qquad\operatorname{ord}_t(T'_1)=b-1.
\label{eq:T-orders}\end{equation}

Let $c$ be the minimum valuation of the entries of $T'$. Then $c\leq a-1$,
and $\widehat T=t^{-c}T'$ is regular and its evaluation at $t=0$ is nonzero. Moreover
\begin{equation}\operatorname{ord}_t(\widehat T_1)\geq b-a>0.
 \label{eq:T1-vanish}\end{equation}

Since $\gamma$'s generic point lies in $Y_{\mathrm{sm}}$, and
normalization is an isomorphism over the smooth locus,
\AIadd{by the finite birational normal-target criterion\cite[Tag~0AB1]{Stacks},} $\widetilde Y$ and
$Y$ agree there; working over $\mathbb C((t))$, we may therefore regard
$v$ as a local parametrization of $Y$ itself near $\gamma$'s
generic point.
Since $(t,z_2,\ldots,z_m)$ is a full coordinate system on the
$m$-dimensional $\widetilde Y$, the $m$ vectors $T,V_2,\ldots,V_m$ form a
tangent basis at that point; together with $v$ itself, they span the
affine cone over the projective tangent $m$-plane to $Y$ along $\gamma$.
Since $T'$ differs from $T$ only by a combination of $V_2,\ldots,V_m$, and
$\widehat T$ differs from $T'$ only by multiplication by the invertible
scalar $t^{-c}\in\mathbb C((t))^\times$, replacing $T$ by $T'$ and then by
$\widehat T$ does not change this span.

Also evaluated at $t = 0$, the columns $v,V_2,\ldots,V_m,\widehat T$ are independent vectors. Indeed,
$v_j=1$, while the other vectors evaluated at $t=0$ have $j$th coordinate zero. $V_\alpha(0)$ were chosen to be independent among themselves, and by construction $\widehat T(0)\neq 0$. \AIreplace{If $\widehat T(0)$ were a combination of $v(0),V_2(0),\ldots,V_m(0)$, matching the zero entries in rows $i_1,\ldots,i_{m-1}$ would
force the $V_\alpha(0)$-coefficients to vanish (by invertibility of the
minor), and matching row $j$ would then force the $v(0)$-coefficient to
vanish too (since $v_j(0)=1\ne0=\widehat T_j(0)$), contradicting that
$\widehat T(0)$ is nonzero, hence also $\widehat T(0)$ is independent from the other vectors.}{If $\widehat T(0)$ were a combination of $v(0),V_2(0),\ldots,V_m(0)$, matching row $j$ first forces the $v(0)$-coefficient to vanish (since $v_j(0)=1$ while all other vectors have $j$th entry $0$); then matching the rows $i_1,\ldots,i_{m-1}$, where $\widehat T(0)$ vanishes, forces the $V_\alpha(0)$-coefficients to vanish by invertibility of the minor. So $\widehat T(0)=0$, a contradiction.}
$v, \widehat T, V_2,\dots, V_m$ therefore span the direct-summand tangent
lattice $L$ required by Lemma~\ref{lem:tangent-limit} and their evaluations at $t = 0$ span the reduction $L_0$.
As a remark, we note that $L_0$ is the Grassmannian limit.
\AIcomment{GEO-02}{\textbf{Proposed change --- author acceptance pending.} In the independence argument above, use row $j$
first: it forces the coefficient of $v(0)$ to be zero. Only then do the
selected minor rows force all $V_\alpha(0)$ coefficients to vanish.
Otherwise the selected rows may contain a contribution from $v(0)$.
Nonzero $\widehat T(0)$ gives the contradiction and a unit maximal minor
gives the direct-summand lattice.}
All vectors $v, \widehat T, V_2,\dots, V_m$ evaluated at $t=0$
have zero $Z_1$-coordinate, so
\[ L_0\subset\{Z_1=0\}.\]
By Lemma~\ref{lem:tangent-limit} $[dx_1] \in \operatorname{Dir}(Y).$
\end{proof}
\end{comment}


\AIadd{\formalref[{any algebraically closed field of characteristic $0$; (i) is the affine-cone form used in Section~\ref{sec:involutive-review}}]{%
\leandecl{Stafford38/Geometry/GeneralAsymptoticConormal.lean}{42}{Stafford38.Geometry.GeneralAsymptoticConormal.coordinate_axis_mem_projective_conormal_directions};
\leandecl{Stafford38/Geometry/GeneralDivisorialVisibleFrame.lean}{55}{Stafford38.Geometry.GeneralDivisorialVisibleFrame.generalDivisorialVisibleFrameWithResidueAlgebraicity};
\leandecl{Stafford38/Geometry/PaperSameWitnessTangentDimension.lean}{70}{Stafford38.Geometry.PaperSameWitnessTangentDimension.tangent_finrank_le_residue_add_one_of_witness};
\leandecl{Stafford38/Geometry/PaperActualWitnessConormalData.lean}{59}{Stafford38.Geometry.PaperActualWitnessConormalData.exists_regularizedOneRowConormalData_of_actual_witness};
\leandecl{Stafford38/Geometry/GeneralConormalAxis.lean}{53}{Stafford38.Geometry.GeneralConormalAxis.exists_groundConormalAxis_of_minimalPrime_unit_transcendental};
\leandecl{Stafford38/Geometry/GeneralAsymptoticLaurentAxis.lean}{31}{Stafford38.Geometry.GeneralAsymptoticLaurentAxis.exists_groundConormalAxis_of_prime_coordinate_avoidance};
\leandecl{Stafford38/Geometry/GeneralAsymptoticConormal.lean}{32}{Stafford38.Geometry.GeneralAsymptoticConormal.coordinate_axis_mem_smooth_fibre_closure};
\leandecl{Stafford38/Geometry/GeneralAsymptoticConormal.lean}{42}{Stafford38.Geometry.GeneralAsymptoticConormal.coordinate_axis_mem_projective_conormal_directions}}}

\section{The involutive obstruction}\label{sec:involutive-review}
\begin{definition}[Fibre-conical]\label{def:fibre-can}
A closed subset $W\subseteq T^*X$ is \emph{fibre-conical} if it is stable under the fibrewise scaling $(y,\xi)\mapsto(y,t\xi)$, $t\in\mathbb{G}_m$.
\end{definition}

%\begin{theorem}[Involutive exclusion]\label{thm:coisotropic-exclusion} %coistropic is equivalent with involutive but involutive is what is actually used
%Let $W\subset T^*X$ be a reduced closed fibre-conical and involutive subset. Let
%$P\in K[\xi_1,\ldots,\xi_n]$ be homogeneous and suppose
%$P(1,0,\ldots,0)\neq0$.
%There is no nonempty $W$ satisfying
%\[
% W\cap\pi^{-1}(H)=\varnothing,\qquad W\subset V(P).
%\]
%\end{theorem}

\begin{theorem}[Hilbert's Nullstellensatz]\label{thm:nss}
Let $K$ be an algebraically closed field and $J\subseteq K[z_1,\dots,z_N]$ an
ideal. Then 
\begin{enumerate}[label=(\alph*)]
    \item $V(J)=\varnothing\iff 1 \in J$
    \item $I(V(J))=\sqrt J$
\end{enumerate}
\end{theorem}

\begin{theorem}[Involutive exclusion]\label{thm:coisotropic-exclusion}
Let $K$ be an algebraically closed field of characteristic zero, and let
$W\subseteq T^*X$ be a nonempty Zariski closed subset that is fibre-conical and
involutive. Let $P\in K[\xi_1,\dots,\xi_n]$ be
a polynomial with $P(1,0,\dots,0)\neq0$. If $P$ vanishes on $W$, then
$W\cap\pi^{-1}(H)\neq\varnothing$.
\end{theorem}

\AIadd{\formalref[{any algebraically closed field of characteristic $0$; $P$ need not be homogeneous}]{%
\leandecl{Stafford38/Geometry/GeneralCoisotropicSets.lean}{160}{Stafford38.Geometry.GeneralCoisotropicSets.exists_zero_base_coordinate_of_isFibreConical};
\leandecl{Stafford38/Geometry/GeneralCoisotropicExclusion.lean}{84}{Stafford38.Geometry.GeneralCoisotropicExclusion.exists_zero_base_coordinate};
\leandecl{Stafford38/Characteristic/BaseRelativePoisson.lean}{29}{Stafford38.Characteristic.BaseRelativePoisson.IsBaseRelativePoisson};
\leandecl{Stafford38/Geometry/GeneralCoisotropicSetsTest.lean}{104}{Stafford38.Geometry.GeneralCoisotropicSetsTest.exact_complex_manuscript_coisotropic_consumer}}}

\begin{proof}
Suppose we have a nonempty closed subset $W\subseteq T^*X$ that is fibre-conical and involutive. Suppose further, that $W\cap\pi^{-1}(H)=\varnothing$. 
We define the (radical) vanishing ideal $J=I(W)\subseteq K[x,\xi]$ consisting of all polynomials that vanish at every point $(y,\xi)\in W$, and $B=J\cap K[x]$, consisting of the functions $f(x)$ with $f(y)=0$ for every $(y,\xi) \in W$. \\

Because $W$ is fibre-conical, the $\xi$-homogeneous components of any
polynomial in $J$ again lie in $J$: if $g=\sum_dg_d\in J$ and
$(y,\xi)\in W$, then $\sum_da^dg_d(y,\xi)=g(y,a\xi)=0$ for all $a\in K^\times$,
and a polynomial in $a$ with infinitely many roots is zero, so every
$g_d(y,\xi)=0$. Hence $J$ is $\xi$-homogeneous.\\

Let $g \in J$ be written in its $\xi$-homogeneous components form, that all lie in $J$. In particular, also $g_0 \in J$. Since $g_0$ has $\xi$-degree $0$, it is a polynomial in $x$ alone, so $g_0 \in B$. Let $Y = V(B)$. For $y \in Y$ we therefore have $g(y, 0) = g_0(y) = 0$, so every $g \in J$ vanishes at $(y,0)$, i.e. $(y,0)\in V(J)$. Since $W$ is closed, $V(J) = V(I(W)) = W$, hence $(y, 0) \in W$, or $Y \times \{0\} \subseteq W$. 
Furthermore, $Y$ is the projection of $W$ and is therefore per definition of $Y$ closed: $\pi(W)\subseteq Y$ since for ($y',\xi')\in W$ and $f \in B\subseteq J$, we have $f(y') = f(y',\xi') = 0$ so $y' \in V(B) = Y$, and $Y\subseteq \pi(W)$, since if $y \in Y$, then also $(y,0)\in W$, so $y = \pi(y,0) \in \pi(W)$.\\
Lastly, assume there exists a $y \in Y\cap H$, then $(y,0)\in W$ and $\pi(y,0) = y \in H$, which would put $(y,0)\in W\cap \pi^{-1}(H)$, which is wrong by our assumption. Therefore, $Y \cap H = \varnothing$.\\

Let $f\in B$. Then $f\circ\pi\in J$.
The Hamiltonian derivation of $f$ applied to a $a \in J$ is
\[
 H_f g=\sum_i\frac{\partial f}{\partial x_i}
 \frac{\partial g}{\partial\xi_i} = \{f,g\}
\]
since the second term of the Poisson bracket vanishes because $f$ does not depend on $\xi$. Since $f, g \in J$ and $\{J, J\} \subseteq J$ by the involutivity of $W$, it follows that $H_f g \in J$ and by induction also $H^m_fg \in J$. 
For a fixed $y \in Y$ and $g \in J$, we define the polynomial $\varphi(t)=g(y,t\,df_y)\in K[t]$. Its derivative is $\varphi'(t) = \sum_i\partial_{x_i}f(y)\frac{\partial g}{\partial \xi_i}{(y, tdf_y)} = H_fg(y, tdf_y)$, so by iteration, $\varphi^{(m)}(t) = H^m_fg(y, tdf_y)$ for $m\geq 0.$ 
Evaluating at $t = 0$, $\varphi^{(m)}(0) = H_f^mg(y,0).$ Since we have established that $(y,0)\in W$ and $H^m_fg \in J$, and all polynomials in $J$ vanish on $W$, we have $\varphi^{(m)}(0) = 0$ for all $m\geq 0$ -- so $\varphi(t)$ is a polynomial it $t$ all of whose derivatives vanish at $t = 0$.  If we write $\varphi(t) = \sum_m c_mt^m$, then $\varphi^{(m)}(0) = m!c_m$ and since $m!\neq 0$ because $K$ has characteristic zero, $c_m = 0$ for all $m$, and therefore, $\varphi(t) = 0$ everywhere. Remembering the definition of $\phi = g(y,t\,df_y) = 0$, so $(y,t\,df_y) \in V(J) = W$.
More generally, since any linear combination of elements in $B$ is again in $B$, we can apply this result to $f = \sum_a c_af_a \in B$ and evaluating at $t = 1$, it follows that 
\[(y, df_y) = (y, d(\sum_ac_af_a)_y) = (y, \sum_a c_ad(f_a)_y) \in W,\]
so the whole linear span $\{y\}\times \operatorname{span}\{df_y:f \in B\}$ is a subset of $W$.\\


Since $V(B)=Y=\pi(W)\neq\varnothing$, $B \neq K[x]$, so $B$ must have at least one minimal prime. The ideal $B$ is radical (by construction since it is a vanishing ideal) in the Noetherian ring $K[x]$, so $B$ has only finitely many minimal primes $\mathfrak q_1,\dots,\mathfrak q_r$ and $B=\bigcap_i\mathfrak q_i$. We choose $\mathfrak q=\mathfrak q_1$ and set $Y_0=V(\mathfrak q)$. Since $\mathfrak q$ is prime, $1\notin\mathfrak q$, so $Y_0\neq\varnothing$ by the Nullstellensatz (Theorem~\ref{thm:nss}); moreover $Y_0$ is closed and irreducible since $\mathfrak q$ is prime \cite[Cor.~I.1.4]{hartshorne}, and $Y_0\subseteq Y$, hence
$Y_0\cap H=\varnothing$, as we have already established $Y\cap H=\varnothing$.

Furthermore, if $r = 1$, we choose $s = 1$, then $sh \in \mathfrak q = B$, and if $r\geq 2$, choose for every $i\geq 2$ elemens $s_i \in \mathfrak q_i \setminus \mathfrak q$. Then $s = s_2s_3\dots s_r$ does not lie in $\mathfrak q$ since no factor lies in $\mathfrak q,$ but does lie in every other $\mathfrak q_i$ since they are ideals and one of the factors lies in them. On the other hand, if $h \in \mathfrak q$, then $sh \in \mathfrak q$ since $h\in \mathfrak q$, and $sh \in \mathfrak q_i$ since $s \in \mathfrak q_i$ for all $i\geq 2$. Therefore, $sh \in \bigcap_i \mathfrak q_i = B,$ i.e. $s\mathfrak q \subseteq B.$ We now fix such a $h \in \mathfrak q,$ and define $Y_0 = V(\mathfrak q).$ For a $y \in Y_0\cap\{s\neq0\}$, $(y, dh_y) \in W:$
Since $sh \in B,$ and $d(sh)_y = s(y)dh_y+h(y)ds_y = s(y)dh_y$ since $h(y) = 0$, so $(y, s(y)dh_y) \in W$ and since $s(y)\neq 0,$ also $(y, dh_y)\in W$. 
For $g\in J$, the polynomial $G(x)=g(x,dh_x)$ therefore vanishes on $Y_0\cap\{s\neq0\}$, and $sG$ vanishes on all of $Y_0$, so $sG \in \mathfrak q$ and therefore $G \in \mathfrak q$, which means that $G$ vanishes on all of $Y_0$. Since this is true for every $g \in J$, we have $g(y, dh_y = G(y) = 0$ for all $g\in J$, so $(y,dh_y) \in V(J)= W.$  As before, we can argue that since $\mathfrak q$ is an ideal, it is closed under $K$-linear combinations, so for $h_1,\dots ,h_k \in \mathfrak q$ and $c_a \in K$, also $(y, \sum_a c_a (dh_a)_y) \in W$, or put differently, $\{y\} \times \operatorname{span}\{dh_y: h \in \mathfrak q\} \subseteq W$ for all $y \in Y_0$. At a smooth point $y$ of $Y_0$, by Definition \ref{def:conormal}, the conormal space is 
\[N_y^*Y_0 = \operatorname{span}\{df_y : f \in \mathfrak q\}\]
since $\mathfrak q$ is prime and hence radical, and so $\mathfrak q = I(Y_0).$ We therefore have that $N_y^*Y_0 \subseteq \{y\} \times W$ for every smooth $y \in Y_0$, and hence $N(Y_0) \subseteq W.$

By assumption, $P$ vanishes on $W$ and hence on $N(Y_0)$. Because we have established, that $Y_0$ is closed, irreducible, and $Y_0 \cap H = \varnothing $, $Y_0$ therefore fulfills all requirements of Theorem \ref{thm:asymptotic-conormal-formal}, so $P(1,0,\dots,0) = 0,$ contradicting our initial assumption. Therefore, $W \cap \pi^{-1}(H) \neq \varnothing.$
\end{proof}


\AIadd{\formalref{%
\leandecl{Stafford38/Geometry/GeneralCoisotropicExclusion.lean}{84}{Stafford38.Geometry.GeneralCoisotropicExclusion.exists_zero_base_coordinate};
\leandecl{Stafford38/Characteristic/BaseRelativePoisson.lean}{72}{Stafford38.Characteristic.BaseRelativePoisson.zeroSection_stable_under_differential_translation_of_isBaseRelativePoisson};
\leandecl{Stafford38/Geometry/GeneralComponentConormalContainment.lean}{72}{Stafford38.Geometry.GeneralComponentConormalContainment.equationConormalClosure_minimalPrime_subset_zeroLocus};
\leandecl{Stafford38/Geometry/GeneralAsymptoticConormal.lean}{32}{Stafford38.Geometry.GeneralAsymptoticConormal.coordinate_axis_mem_smooth_fibre_closure}}}


\section{Proof of Theorem~\ref{thm:main}}\label{sec:proof}

\begin{lemma}\label{lem:base-change}
Let $K/k$ be a field extension and put
$I_K=\tilde dA_n(K)+x_1^N\tilde dA_n(K)$, $Q_K=Q\otimes_kK=A_n(K)/I_K$, with the
order filtration $Q_{K,i}=q\,F_iA_n(K)$. Then
\[
 \gr Q_K\cong\gr Q\otimes_kK,\qquad
 \Ann_{K[x,\xi]}\gr Q_K=\bigl(\Ann_{k[x,\xi]}\gr Q\bigr)\,K[x,\xi],
\]
and $Q=0\iff Q_K=0$. In particular, $\Char(Q_K)$ is the preimage of $\Char(Q)$
under $\Spec K[x,\xi]\to\Spec k[x,\xi]$.
\end{lemma}
\begin{proof}
By the PBW basis, $F_iA_n(K)=F_iA_n(k)\otimes_kK$ and $I_K=I\otimes_kK$, so
$Q_{K,i}=Q_i\otimes_kK$. As $K$ is flat over $k$,
$Q_{K,i}/Q_{K,i-1}=(Q_i/Q_{i-1})\otimes_kK$, and summing over $i$ gives
$\gr Q_K=\gr Q\otimes_kK$. The module $\gr Q$ is finitely generated (cyclic,
generated by the class of $q$) over the Noetherian ring $k[x,\xi]$, and
$k[x,\xi]\to K[x,\xi]$ is flat, so the annihilator of the finite module
commutes with this base change \cite[Tag~05DT]{Stacks}. Finally, $Q\otimes_kK=0$
implies $Q=0$, because $K$ is a nonzero free $k$-module.
\end{proof}

\noindent\textbf{Proof of Theorem~\ref{thm:main}.}
Let $d\in A_n(k)$ be nonzero of Bernstein degree $N$ and let
$\tilde d=c^{-1}\Phi_g(d)$ be its monic normal form
(Remark~\ref{rem:transport}). If $n=0$, then $A_0(k)=k$ and every nonzero $d$ is
a unit anyways, so let $n\ge1$. As in Section \ref{sec:euler} we define
\[ I=\tilde dA_n(k)+x_1^N\tilde dA_n(k),\qquad Q=A_n(k)/I .\]
It suffices to show $Q=0$: then $1\in I$, i.e.\ there are
$R,S\in A_n(k)$ with $1=\tilde dR+x_1^N\tilde dS$.

Let $K=\overline k$, and let
\[W=\Char(Q_K)=V\bigl(\Ann_{K[x,\xi]}\gr Q_K\bigr)\subseteq K^n\times K^n,\]
with vanishing ideal $I(W)=\sqrt{\Ann\gr Q_K}$ by the Nullstellensatz (Theorem \ref{thm:nss}). We check
the hypotheses of Theorem~\ref{thm:coisotropic-exclusion} for $W$ and the
polynomial $P$ of Section~\ref{sec:char}:

\begin{itemize}
\item $W$ is Zariski closed, since it is the zero set of an ideal of polynomials.
\item The module $\gr Q_K$ is graded by the order filtration, which assigns degree
$0$ to the $x_i$ and degree $1$ to the $\xi_i$; so $\Ann\gr Q_K$  is homogeneous in $\xi$. Hence if $(y,\xi)\in W$
and $g\in\Ann\gr Q_K$, then $g(y,t\xi)=\sum_dt^dg_d(y,\xi)=0$ for every
$t\in K^\times$, where $g_d$ denotes the $\xi$-homogeneous component of degree $d$ of $g$. So
$(y,t\xi)\in W$, i.e.\ $W$ is fibre-conical.
\item By Gabber's theorem (Theorem~\ref{ithm:gabber}) $I(W)$ is involutive. 
\item By eq.~\eqref{eq:P-annihilates}, $P\in
\Ann_{k[x,\xi]}\gr Q$, hence $P\in\Ann_{K[x,\xi]}\gr Q_K$ by
Lemma~\ref{lem:base-change}, so $P$ vanishes on $W$.
\item By eq.~\eqref{eq:P-e1}, $P(1,0,\dots,0)=1\neq 0$.
\item  $W=\Char(Q_K)$ is the preimage of $\Char(Q)$ under
$\Spec K[x,\xi]\to\Spec k[x,\xi]$ (Lemma~\ref{lem:base-change}), and
$\pi^{-1}(H)=V(x_1)$ over $K$ is the preimage of $\pi^{-1}(H)=V(x_1)$ over $k$.
By Proposition~\ref{prop:char-avoids}, $\Char(Q)\cap\pi^{-1}(H)=\varnothing$, and
the preimage of an empty set is also empty, therefore $W\cap\pi^{-1}(H)=\varnothing$.

\end{itemize}
By Theorem~\ref{thm:coisotropic-exclusion}, $W=\varnothing$. As $K$ is
algebraically closed, by the Nullstellensatz (Theorem \ref{thm:nss})  $1 \in \Ann_{K[x,\xi]}\gr Q_K$,
so $\gr Q_K=0$. Since the filtration of $Q_K$ is exhaustive with $Q_{K,-1}=0$,
induction on $m$ implies $Q_{K,m}=0$ for all $m$, so $Q_K=0$, and
$Q=0$ by Lemma~\ref{lem:base-change}. Therefore, $A_n(k) = \tilde dA_n(k) + x_1^N\tilde dA_n(k)$, or equivalently $1=\tilde dR+x_1^N\tilde dS$
with $R,S\in A_n(k)$.

Lastly, we apply the inverse of the automorphism $\Phi_g$ of
Remark~\ref{rem:transport}. Since $\tilde d=c^{-1}\Phi_g(d)$ with $c\in k^\times$
central, this gives
\[
 1=d\bigl(c^{-1}\Phi_g^{-1}(R)\bigr)+F\,d\bigl(c^{-1}\Phi_g^{-1}(S)\bigr),
 \qquad F=\Phi_g^{-1}(x_1)^N,
\]
which proves Theorem~\ref{thm:main}. \qed



\AIadd{\formalref{\leandecl{Stafford38/FoundationClosure.lean}{41}{Stafford38.FoundationClosure.canonicalSupportVanishingViaGeneralCoisotropic}:
exclusion over $\overline k$ \leandecl{Stafford38/Geometry/GeneralCoisotropicCanonicalAdapter.lean}{222}{Stafford38.Geometry.GeneralCoisotropicCanonicalAdapter.algebraicallyClosedCanonicalSupportVanishing_of_generalCoisotropic},
descent \leandecl{Stafford38/Weyl/FilteredScalarLifting.lean}{472}{Stafford38.Weyl.FilteredScalarLifting.canonicalSupportDescent},
certificate \leandecl{Stafford38/Characteristic/CanonicalCertificate.lean}{33}{Stafford38.CharacteristicCanonicalCertificate.exists_fixedSource_certificate_of_orderCharacteristicSupport_eq_empty}}}

\AIadd{\formalref[{stated for the canonical ideal; the argument uses nothing else}]{%
\leandecl{Stafford38/Weyl/FilteredScalarLifting.lean}{449}{Stafford38.Weyl.FilteredScalarLifting.presentedWeylScalarExtension_map_orderInitialIdeal_eq};
\leandecl{Stafford38/Weyl/FilteredScalarLifting.lean}{387}{Stafford38.Weyl.FilteredScalarLifting.target_orderInitialIdeal_le_map_source_orderInitialIdeal};
\leandecl{Stafford38/Weyl/PresentedScalarExtension.lean}{390}{Stafford38.Weyl.PresentedScalarExtension.presentedWeylScalarExtension_map_orderInitialIdeal_le};
\leandecl{Stafford38/Characteristic/EmptySupportVanishing.lean}{83}{Stafford38.CharacteristicEmptySupportVanishing.orderInitialIdeal_eq_top_iff_rightQuotient_subsingleton};
\leandecl{Stafford38/Characteristic/GeometricSupportDescent.lean}{101}{Stafford38.Characteristic.GeometricSupportDescent.map_scalarPolynomialMap_eq_top_iff}}}


Lastly, as Stafford already showed, from $1 = dR+FdS$ follows cyclicity of the torsion right module:
\begin{corollary}\label{cor:cyclic}
Every finitely generated torsion right $A_n(k)$-module is cyclic.
\end{corollary}

\begin{proof}
    Suppose for every $d\neq 0 \in A_n$, there exists $F, R, S \in A_n$, such that $1 = dR + FdS$. Let us denote a torsion module by $T$. Stafford's theorem (Theorem 3.7 \mbox{\cite[p.~438]{Stafford1978}}) established, that every finitely generated torsion module can be generated by 2 generators, so let us denote the generators of $T$ with $t_1$ and $t_2$. We can define a homomorphism
    \[\phi:A_n \oplus A_n\rightarrow T \quad \phi(r,s)\mapsto t_1r+t_2s.\]
    This map is surjective, since $t_1, t_2$ generate all of $T$. By the first isomorphism theorem (for $\phi:M\rightarrow N$, $im(\phi) \cong M/ker(\phi)$), it follows
    \[T \cong (A_n\oplus A_n)/M\]
    with $M = ker(\phi)$. \\
    $T$ is torsion and therefore $\forall t\in T: \exists m \in A_n\backslash\{0\}: tm = 0.$ Specifically, let us denote with $d_1$ and $d_2$ the annihilators of $t_1$ and $t_2$, i.e.
    \[t_1d_1 = t_2d_2 = 0.\]
    $A_n$ is a right Ore domain, which means it has the property that for any nonzero $d_1, d_2$ there exists elements $r_1, r_2 \in A_n$ such that $d_1r_1 = d_2r_2 \neq 0.$
    Defining therefore $d := d_1r_1 = d_2r_2$ yields us a common annihilator of $t_1$ and $t_2$.\\
    $(d, 0)$ and $(0,d)$ are elements in $M = ker(\phi)$, since $\phi(d, 0) = t_1d = 0$ and $\phi(0, d) = t_2d = 0.$\\
    Because $(0,d), (d, 0) \in M$, also
    $(d, 0)S + (0, d)R = (dS, dR) = (dS, 1-FdS) =: (a, 1-Fa) \in M$, i.e. $(a, 1-Fa)$ is an element of $M$. \\
    On the other hand, $A_n\oplus A_n$ can be generated by $(a, 1-Fa)$ and $(1, -F)$. We can check this by finding $u,v$ such that for all $(x,y)\in A_n\oplus A_n$ we can write
    \[(x,y) = (a, 1-Fa)u + (1, -F)v.\]
    Suitable choices are $u = y+Fx$ and $v = x-a(y+Fx)$.\\
    Lastly, since  $A_n\oplus A_n$ can be generated by $(a, 1-Fa)$ and $(1, -F)$ and $(a, 1-Fa)$ is an element of $M$, it follows that $T \cong (A_n\oplus A_n)/M$ can be generated by the image of $(1, -F)$, i.e. $T$ is cyclic.
\end{proof}
\AIadd{\formalref{\leandecl{Stafford38/TorsionCyclicity.lean}{66}{Stafford38.TorsionCyclicity.weyl_isCyclic_of_isRightTorsion},
paper generator \leandecl{Stafford38/TorsionCyclicity.lean}{39}{Stafford38.TorsionCyclicity.exists_span_adjusted_pair},
two generators \leandecl{Stafford38/TorsionCyclicity.lean}{52}{Stafford38.TorsionCyclicity.exists_span_singleton_eq_span_pair},
domain \leandecl{Stafford38/Weyl/Domain.lean}{25}{Stafford38.WeylDomain.mul_ne_zero};
Mathlib-only statement \leandecl{CorollaryChallenge.lean}{34}{Stafford38CorollaryChallenge.torsionCyclicStatement}}}

\section{Outlook}
\AIremove{We have proved}
\AIadd{The cofactors $R,S$ are obtained non-constructively: the argument proves
$Q=0$ without producing a certificate. Effective bounds for $\deg_BR$ and
$\deg_BS$ in terms of $\deg_Bd$, and an algorithm in the spirit of the
constructive two-generator methods, remain open. The Global Stafford
extension raises the same questions for rings of differential operators on
smooth affine varieties.}



\clearpage
\bibliographystyle{amsplain}
\bibliography{references}






\begin{appendix}


% Evidence: project ledgers, archived outward session messages, model-use
% records, independent reviews, and canonical formal-proof provenance.
\section{AI-assisted research and formalization}\label{app:ai-workflow}

Proof discovery, counterexample searches, review, and formalization were
assisted by agentic AI tools. The work began with
cyclic-extension computations in Macaulay2 and Maple. Attempts to construct
universal cofactors directly produced special cases and reductions but
eventually stalled. The decisive change was to study the canonical quotient
$Q=A_n/(dA_n+x^N dA_n)$: Euler surjectivity and monicity constrained its
characteristic support, while Gabber's involutivity and asymptotic conormal
geometry forced that support to be empty. This turned cofactor selection
into a quotient-vanishing problem. Full Lean formalization was deferred
until the paper argument was complete, although small formal checks
continued throughout. In the final formal development, support avoidance
over the coordinate hyperplane was proved by a localized filtered Koszul
length contradiction.

The workflow separated proposed ideas from verified claims. One controller
maintained the authoritative proof state, while parallel workers explored
distinct mechanisms and returned evidence with an explicit scope.
Independent reviewers checked fixed revisions for quantifiers, sidedness,
circularity, and missing hypotheses. The project protocols were based
on a \texttt{math-frontier} agent skill developed for this purpose: freeze the
claim, identify its first unresolved dependency, and specify how a route
could fail. Conceptual prompts requested a small number of distinct
mechanisms; paper-only tasks sought complete arguments before implementation.
A condensed representative prompt, adapted from the archived protocols, is:
\begin{quote}\small
Work on one precisely stated unresolved claim. Propose at most three
distinct mechanisms, each with an intermediate lemma, a noncircular
justification, and an adversarial test. Preserve the quantifiers and
right-sided cofactors. For a paper task, supply a complete argument; for
formalization, prove the required hypotheses without adding axioms or
using \texttt{sorry}. Return evidence and the first unsupported bridge; do not
promote a conditional result to a theorem.
\end{quote}

Research and review sessions used GPT-5.5, GPT-5.6 Luna and Sol,
GPT-6 Astra, and Claude Fable 5, Haiku 4.5 and Opus 5. The project protocols
assigned routine implementation and replay to Luna and conceptual
consultation to Sol; independent reports also record Opus reviews.
Subsequent manuscript, proof-map, build and release work used Fable 5.1,
Opus 5, Sonnet 5 and GPT-6 Luna. These records establish participation and task roles,
without assigning a particular mathematical idea to one model.

The completed proof uses Lean~4 and Mathlib and pins AlgebraicAnalysis v0.3.0
\cite{Lean4,Mathlib2020,AlgebraicAnalysis030}; the later library archive is v0.3.3~\cite{AlgebraicAnalysis033}.
Formalization distinguished theorems proving required inputs from
conditional theorems merely consuming them. Statement comparison and
independent NanoDa and Lean kernel checks verified the submitted declarations
\cite{Comparator,NanoDa,PalomarDocs}; their proofs use only
\texttt{propext}, \texttt{Classical.choice}, and \texttt{Quot.sound}.
The Stafford38 and Global Stafford headline proofs are registered separately in
Palomar~\cite{Stafford38Palomar2026,GlobalStaffordPalomar2026}; their Lean developments are archived~\cite{Stafford38Zenodo123,GlobalStaffordZenodo106}.

\section{Further work}
\AIcomment{REF-01}{\textbf{Proposed change --- author acceptance pending.} Please confirm that the references here point to the order-symbol result and the characteristic-containment statement used in the argument.}
\AIreplace{AI also proved some related stuff regarding Right Ore localizations, the corresponding left-hand identity, and apparently the corresponding identity over other rings such as differential rings with laurent series coefficients; I am leaving that out for now.}{The formalization also proves the following consequences, which are not otherwise used here:
the identity in every right Ore localization of $A_n(k)$
\formalref{\leandecl{Stafford38/LocalizationCorollaries.lean}{31}{Stafford38.LocalizationCorollaries.s38_rightOreLocalization}};
the left-handed identity $1=Rd+SdF$ via the formal adjoint
\formalref{\leandecl{Stafford38/LeftHandedCorollary.lean}{23}{Stafford38.LeftHandedCorollary.leftHanded_of_universalStatement}};
the identity for finite-order differential operators on localizations of $k[x_1,\dots,x_n]$, including principal opens, Laurent-polynomial coefficients and the fraction field
\formalref{\leandecl{Stafford38/LocalizedDifferentialCorollaries.lean}{83}{Stafford38.LocalizedDifferentialCorollaries.s38_unconditional_localized_differential}};
and, independently of the geometric argument, an evolutionary version with arbitrary coefficients commuting with one Weyl pair
\formalref{\leandecl{Stafford38/EvolutionaryCorollary.lean}{36}{Stafford38.Evolution.evolutionaryCorollary}}.}
\AIcomment{SCOPE-01}{\textbf{Proposed change --- author acceptance pending.} Review the scope before accepting this paragraph. The coefficients in the stated localization result are Laurent polynomials, not arbitrary Laurent series, and the operators have finite order. Retain further consequences only with their exact hypotheses and proof references; optional formal declaration links are available in the technical supplement.}

\end{appendix}


\end{document}
